% =====================================================================
%  The Frankl cohomology, re-founded on Fork A.
%  Phase 1: the LaTeX re-foundation.
%
%  Work item mg-c15b (Frankl-ForkA-Refound-LaTeX).
%  Per the Daniel 2026-05-22 directive "frankl - commit": commit to the
%  Fork A re-foundation, the route the vision-check (mg-6edc) and the
%  n=3 Fork-D probe (mg-565a) resolved to.
%
%  REVISED under work item mg-365e (Frankl-ForkA-Phase2-Revision) per
%  the mg-d300 independent Phase-2 validation (AMBER): two surviving
%  over-claims corrected by honest re-statement -- the Y1 cofiber-LES
%  reduction is only Stab(x) = S_{n-1} equivariant (not S_n), so R1
%  acquires an S_{n-1}->S_n branching sub-residual; and the residual
%  set is not "exactly three" -- R3 is a three-part edge-map debt
%  (existence, degree, injectivity). See Status~\ref{status:mg365e-revision}.
%
%  REVISED AGAIN under work item mg-894c (Frankl-ForkA-Phase2-Revision2)
%  per the mg-5221 independent Phase-2 re-validation (AMBER): a third
%  over-claim of the identical shape -- the twisted-bridge null-homotopy
%  iota_x* ~ 0 presented as a discharged ported input -- demoted to the
%  open sub-residual R1c; plus an exhaustive sweep of every ported
%  Walsh/chi_S fact, certifying there is no fourth prong. See
%  Status~\ref{status:mg894c-revision} and the ledger
%  docs/Frankl-ForkA-Phase2-Revision2.md.
%
%  This is a STANDALONE document. It is Phase 1 (math-before-Lean) of a
%  multi-phase arc: Phase 1 LaTeX (this document), Phase 2 independent
%  validation, Phase 3 Lean. It does not modify the consolidated paper
%  (../main.tex), which honestly records the pre-Fork-A state.
%
%  Build:  pdflatex refoundation && pdflatex refoundation
%  (self-contained; no bibliography, no external inputs).
% =====================================================================
\documentclass[11pt,oneside]{amsart}

\usepackage[utf8]{inputenc}
\usepackage[T1]{fontenc}
\usepackage{amsmath,amssymb,amsthm,mathtools}
\usepackage{stmaryrd}
\usepackage{enumitem}
\usepackage{booktabs}
\usepackage{longtable}
\usepackage[hidelinks]{hyperref}
\usepackage{xcolor}
\usepackage[margin=1.1in]{geometry}

\setcounter{secnumdepth}{3}
\setcounter{tocdepth}{2}

% ---- theorem environments ------------------------------------------
\theoremstyle{plain}
\newtheorem{theorem}{Theorem}[section]
\newtheorem{lemma}[theorem]{Lemma}
\newtheorem{proposition}[theorem]{Proposition}
\newtheorem{corollary}[theorem]{Corollary}
\newtheorem{conjecture}[theorem]{Conjecture}

\theoremstyle{definition}
\newtheorem{definition}[theorem]{Definition}
\newtheorem{example}[theorem]{Example}
\newtheorem{construction}[theorem]{Construction}

\theoremstyle{remark}
\newtheorem{remark}[theorem]{Remark}
\newtheorem{notation}[theorem]{Notation}

% Honesty / status callouts (same convention as the consolidated paper).
\newtheoremstyle{honeststyle}{6pt}{6pt}{\normalfont}{}%
  {\bfseries}{.}{.5em}{}
\theoremstyle{honeststyle}
\newtheorem{status}[theorem]{Status}
\newtheorem{openresidual}[theorem]{Open residual}
\newtheorem{debt}[theorem]{Owed debt}

% ---- macros ---------------------------------------------------------
\newcommand{\F}{\mathcal{F}}
\newcommand{\G}{\mathcal{G}}
\newcommand{\Fstar}{\F^{\star}}
\newcommand{\bias}{\beta}
\newcommand{\sgn}{\mathrm{sgn}}
\newcommand{\std}{\mathrm{std}}
\newcommand{\triv}{\mathbf{1}}
\newcommand{\Sym}{\mathfrak{S}}
\newcommand{\Z}{\mathbb{Z}}
\newcommand{\Q}{\mathbb{Q}}
\newcommand{\N}{\mathbb{N}}
\newcommand{\Cn}{\mathcal{C}^{\cap}_n}
\newcommand{\Cnm}{\mathcal{C}^{\cap}_{n-1,x}}
\newcommand{\Xbb}{\mathbb{X}}
\newcommand{\Dx}{\mathcal{D}_x}
\newcommand{\hocolim}{\operatorname*{hocolim}}
\newcommand{\ob}{\mathrm{ob}}
\newcommand{\tr}{\mathrm{tr}}
\newcommand{\BK}{\mathrm{BK}}
\newcommand{\Tot}{\mathrm{Tot}}
\newcommand{\fobs}{\mathcal{F}_{\mathrm{obs}}}
\newcommand{\cofib}{\mathrm{cofib}}
\newcommand{\walsh}[1]{\chi_{#1}}
\newcommand{\Ch}{\mathrm{Ch}}
\renewcommand{\thefootnote}{\fnsymbol{footnote}}
\DeclareMathOperator{\im}{im}
\DeclareMathOperator{\Hom}{Hom}

\title[The Frankl cohomology re-founded on Fork A]{%
  The Frankl cohomology, re-founded on Fork A\\
  \normalsize Phase 1: the LaTeX re-foundation}

\author{The \texttt{union\_closed} development}
\date{Work items \texttt{mg-c15b} (Phase 1), \texttt{mg-365e}
  (revision per the \texttt{mg-d300} Phase-2 validation), and
  \texttt{mg-894c} (revision~2 per the \texttt{mg-5221} Phase-2
  re-validation) --- \today\\[2pt]
  \normalfont\small Phase 1 of the Fork A re-foundation arc
  (LaTeX $\to$ independent validation $\to$ Lean)}

\begin{document}

\begin{abstract}
The cohomological arc of the union-closed (Frankl) program was halted
when its Phase-1.5 audit (\texttt{mg-0405}) refuted two load-bearing
objects: the against-trace structure map on which the cohomology object
$X^{\cap}_n$ was built, and the $(\Z/2)^n$-action on which its Walsh
grading rested. A vision-check (\texttt{mg-6edc}) diagnosed both
refutations as consequences of \emph{drift} from the program's stated
goal, and a bounded $n=3$ probe (\texttt{mg-565a}) eliminated the
literal-vision foundation (Fork~D --- the bare nerve of the category,
which is rationally acyclic) while \emph{positively endorsing} Fork~A:
re-found the cohomology as the \emph{covariant} homotopy colimit of the
per-family coefficient diagram $\F\mapsto X(\F)$ over the trace
category $\Cn$, drop the $(\Z/2)^n$-action, grade by $\sgn_{\Sym_n}$,
and keep Walsh only as the harmonic description of the obstruction
cochain. This document is Phase~1 of the resulting re-foundation: the
LaTeX construction.

We carry it out. The covariant homotopy colimit $H^{*}(\Cn;X)$ is a
genuinely defined cohomology theory --- its structure map is the
canonical trace \emph{projection}, which exists and is functorial, so
the Phase-1.5 variance refutation is dodged by construction. Dropping
the $(\Z/2)^n$-action dissolves the second Phase-1.5 refutation
outright: there is no group action left to fail. The grading is
re-established by Maschke's theorem for the genuine $\Sym_n$-action,
with the sphere class identified as the $\sgn_{\Sym_n}$ isotype. We
re-found the obstruction class $\ob(\Fstar)\in H^{n-1}(\Cn;X)$ and the
two-part contradiction. \textbf{We are honest about what the
re-foundation does and does not buy.} It buys a sound, well-defined
framework that clears both Phase-1.5 walls. It does \emph{not} close
Frankl, and Phase~1 was never meant to. The closing contradiction is
re-established as a conditional theorem, and its open inputs are named
precisely and in full: the vanishing statement $Y_1$ (the cofiber
computation --- its target object unchanged from the pre-Fork-A state,
but its reduction now carrying two steps the $(\Z/2)^n\to\Sym_n$
re-grading creates: a genuine $\Sym_{n-1}\to\Sym_n$ branching step,
and a Fork-A-native discharge of the twisted-bridge null-homotopy);
the landing statement $Y_2$ (whose pre-Fork-A mechanism rested on the
now-dropped $(\Z/2)^n$-equivariance and which Fork~A re-states as an
$\Sym_n$ question, localized to the representation theory of the
obstruction sheaf); and a three-part faithfulness debt (whether the
\v{C}ech-to-cohomology edge map that promotes the obstruction class
\emph{exists}, lands in \emph{degree $n-1$}, and is \emph{injective}
--- the vision's clause $V_3$). The first owed debt, the re-grading, is
discharged at the structural level and its remaining open parts (the
$Y_2$ landing, the $\Sym_{n-1}\to\Sym_n$ branching step, and the
twisted-bridge null-homotopy) are precisely localized; the second owed
debt, faithfulness, is precisely localized as the three-part residual.
\textbf{The verdict is GREEN for the Phase-1 scope --- the
re-foundation is delivered and sound --- with the explicit accounting
that Frankl is not proved and that the conditional theorem rests on
the open residuals R1 (three-part: cofiber, branching, and the
twisted-bridge null-homotopy), R2, and R3 (three-part
faithfulness).} This document has been revised twice following
independent Phase-2 validation; see
Status~\ref{status:mg365e-revision} and
Status~\ref{status:mg894c-revision}.
\end{abstract}

\maketitle

\tableofcontents

% =====================================================================
\section*{\S0\quad Verdict (top of page)}
\addcontentsline{toc}{section}{\S0\quad Verdict}
% =====================================================================

\subsection*{0.1\quad Headline}

\noindent\textbf{GREEN --- FORK-A-REFOUNDATION-DELIVERED (Phase~1).}
The Frankl cohomology is re-founded on Fork~A. The new cohomology
object $H^{*}(\Cn;X)$ --- the covariant homotopy colimit of the
per-family cubical-defect diagram over the trace category, with the
canonical projection $\pi_T$ as structure map --- is \emph{genuinely
defined} (Theorem~\ref{thm:welldefined}); this dodges the Phase-1.5
against-trace refutation by construction, because the projection
exists and is functorial where the against-trace map did not exist at
all. The $(\Z/2)^n$-action is \emph{dropped}; this dissolves the
Phase-1.5 $(\Z/2)^n$ refutation outright (Status~\ref{status:c-block}).
The grading is re-established by Maschke for the genuine
$\Sym_n$-action (Theorem~\ref{thm:maschke}). The obstruction class and
the two-part contradiction are re-established on the new object
(\S\ref{sec:obstruction}--\S\ref{sec:contradiction}).

\medskip\noindent
\textbf{Frankl is not proved, and Phase~1 was never meant to prove it.}
The verdict GREEN means: the Phase-1 scope --- construct the object,
re-grade, re-establish the framework, localize the two debts --- is
accomplished, and the re-foundation is internally sound. It does
\emph{not} mean the conjecture is closer to proved than the
pre-Fork-A state recorded. The genuine open content remains, and
Phase~1's task is to name it precisely and in full --- as the
residuals R1 (three-part), R2, and R3 (three-part):
\begin{itemize}[leftmargin=1.6em]
\item[\textbf{R1}] \emph{The $Y_1$ residual} (Open
  residual~\ref{res:y1}), in three parts: \textbf{R1a}, the
  $\Sym_{n-1}$-isotype cohomology of the cofiber of the
  matched-cylinder inclusion --- as a target object unchanged from the
  pre-Fork-A state; \textbf{R1b}, a new $\Sym_{n-1}\to\Sym_n$
  branching step created by the $(\Z/2)^n\to\Sym_n$ re-grading
  (Remark~\ref{rem:y1-branching}); and \textbf{R1c}, the
  twisted-bridge null-homotopy $\iota_x^{*}\simeq0$ --- a
  $\walsh{S}$-graded fact the ported bridge does \emph{not} discharge
  in Fork~A (Remark~\ref{rem:bridge-not-ported}), also created by the
  $(\Z/2)^n\to\Sym_n$ re-grading.
\item[\textbf{R2}] \emph{The $Y_2$ residual} (Open
  residual~\ref{res:y2}): whether the obstruction class avoids the
  $\sgn$ sphere-isotype. The pre-Fork-A $Y_2$ mechanism rested on the
  now-dropped $(\Z/2)^n$-equivariance; Fork~A re-states the landing as
  an $\Sym_n$ question and localizes it.
\item[\textbf{R3}] \emph{Faithfulness} (Owed
  debt~\ref{debt:faithful-localized}), a three-part debt: the
  \emph{existence}, the \emph{degree-$(n-1)$ behaviour}, and the
  \emph{injectivity} of the \v{C}ech-to-cohomology edge map --- the
  vision's clause $V_3$, the wrapper-risk.
\end{itemize}

\subsection*{0.2\quad The Fork A scorecard}

{\small
\begin{longtable}{p{0.28\textwidth}p{0.15\textwidth}p{0.49\textwidth}}
\toprule
\textbf{Component} & \textbf{Status} & \textbf{Where / why}\\
\midrule
\endhead
The covariant cohomology object $H^{*}(\Cn;X)$
 & \textsf{GREEN}
 & Theorem~\ref{thm:welldefined}. The structure map is the projection
   $\pi_T$ (Theorem~\ref{thm:projection}); it exists and is functorial.
   The Phase-1.5 against-trace refutation does not apply.\\
Dropping the $(\Z/2)^n$-action
 & \textsf{GREEN}
 & Status~\ref{status:c-block}. There is no group action left to fail;
   the Phase-1.5 $(\Z/2)^n$ refutation is dissolved, not repaired.\\
$\Sym_n$-grading by isotype
 & \textsf{GREEN}
 & Theorem~\ref{thm:maschke}. Genuine $\Sym_n$-action; Maschke.
   The sphere class is the $\sgn_{\Sym_n}$ isotype
   (\S\ref{ssec:sphere-class}).\\
The obstruction \emph{cocycle} $\{m_{xy}\}$
 & \textsf{GREEN} (as a \v{C}ech cocycle)
 & \S\ref{sec:obstruction}. The mismatch cocycle is constructed
   verbatim from the \v{C}ech double complex of $\fobs$. Its
   \emph{promotion} into $\widetilde H^{n-1}(\Cn;X)$ as $\ob(\Fstar)$
   is the three-part R3 debt (Owed
   debt~\ref{debt:faithful-localized}).\\
$Y_1$-A (the vanishing) --- residual R1
 & \textsf{RED --- open}
 & \S\ref{ssec:y1a}. The cofiber-LES reduction ports to Fork~A as an
   $\Sym_{n-1}$-equivariant statement
   (Theorem~\ref{thm:y1-reduction}); the residual is three-part (Open
   residual~\ref{res:y1}) --- R1a, the cofiber, unchanged as a target
   object; R1b, a new $\Sym_{n-1}\to\Sym_n$ branching step; R1c, the
   twisted-bridge null-homotopy, which the ported bridge does not
   discharge in Fork~A.\\
$Y_2$-A (the landing) --- residual R2
 & \textsf{AMBER --- open}
 & \S\ref{ssec:y2a}. Mechanism re-stated in $\Sym_n$ terms; the
   separation from the sphere class is localized
   (Open residual~\ref{res:y2}).\\
Owed debt 1 --- the re-grading
 & \textsf{ADDRESSED}
 & \S\ref{ssec:debt1}. Structural re-grading discharged
   (Theorem~\ref{thm:maschke}, \S\ref{ssec:sphere-class}); the open
   part is residual R2 \emph{and} the two new $Y_1$-side
   sub-residuals R1b (branching) and R1c (the bridge
   null-homotopy).\\
Owed debt 2 --- faithfulness ($V_3$ / wrapper-risk) --- residual R3
 & \textsf{LOCALIZED}
 & Owed debt~\ref{debt:faithful-localized}. Reduced to one named edge
   map, but a three-part debt on it: \emph{existence},
   \emph{degree-$(n-1)$ landing}, and \emph{injectivity}.\\
\bottomrule
\end{longtable}}

\subsection*{0.3\quad What Phase 1 is, and is not}

Phase~1 is the LaTeX re-foundation: construct the Fork~A object, do the
re-grading, re-establish the obstruction class and the two-part
contradiction, and discharge or precisely localize the two owed debts.
It is \emph{not} a proof of Frankl, and it is not a closure of any of
the residuals R1--R3 (R1 itself three-part, R3 itself three-part). A
RED verdict --- Fork~A hitting a genuine
wall --- was an allowed outcome (the \texttt{mg-c15b} brief sanctions
it); Phase~1 does not return RED, because nothing in the re-foundation
is refuted: the construction is well-defined, the re-grading goes
through, and the contradiction re-assembles as a conditional theorem.
The honest one-sentence summary is \S\ref{ssec:final}. The arc
continues with a further independent re-validation of this
twice-revised document and then Phase~3 (Lean), each gated on the
previous returning GREEN.

\begin{status}[Revision following the \texttt{mg-d300} Phase-2 validation]
\label{status:mg365e-revision}
This document has been revised (work item \texttt{mg-365e}) following
the independent Phase-2 validation \texttt{mg-d300}, which returned
\textsf{AMBER}: the Fork~A core was independently confirmed sound, but
two over-claims were caught that the Phase-1 self-audit had missed.
\begin{enumerate}[label=\textup{(\Alph*)},leftmargin=2.4em]
\item \emph{The $Y_1$-A reduction does not port unchanged.} Theorem~%
  \ref{thm:y1-reduction} was stated as carrying an
  $\Sym_n$-equivariant cohomology long exact sequence. It does not:
  the smaller object $\Xbb_{n-1,x}$ is stable only under the point
  stabiliser $\mathrm{Stab}(x)\cong\Sym_{n-1}$, so the LES is only
  $\Sym_{n-1}$-equivariant. The pre-Fork-A reduction split
  isotype-by-isotype only because the abelian $(\Z/2)^n$ has a
  quotient acting on the smaller cube; $\Sym_n$ has no such quotient.
  The reduction therefore acquires a genuine $\Sym_{n-1}\to\Sym_n$
  branching step, recorded here as the new sub-residual R1b
  (Remark~\ref{rem:y1-branching}, Open residual~\ref{res:y1}).
\item \emph{The residual set is not ``exactly three.''} The
  obstruction class is promoted into $H^{*}(\Cn;X)$ through an edge
  map whose \emph{existence} and \emph{degree-$(n-1)$ behaviour} are
  load-bearing but were not in the formal accounting --- the formerly
  stated R3 named only \emph{injectivity}. R3 is in fact a three-part
  edge-map debt (R3a existence, R3b degree, R3c injectivity; Owed
  debt~\ref{debt:faithful-localized}), and
  Theorem~\ref{thm:contradiction} now lists every input in full.
\end{enumerate}
A minor point --- the proof of Proposition~\ref{prop:sphere-sgn}
assumed more than its stated hypothesis --- is also corrected
(Remark~\ref{rem:sphere-sgn-hyp}). Both findings are honest
re-statements, not refutations: the Fork~A strategy and its core
constructions (Theorems~\ref{thm:projection},
\ref{thm:welldefined}, \ref{thm:maschke}) are unchanged and were
independently confirmed sound by \texttt{mg-d300}. The verdict GREEN
above is for the Phase-1 scope as corrected; an independent
re-validation is the next gate.
\end{status}

\begin{status}[Revision~2 following the \texttt{mg-5221} Phase-2
re-validation]
\label{status:mg894c-revision}
This document has been revised a second time (work item
\texttt{mg-894c}) following the independent Phase-2 re-validation
\texttt{mg-5221}, which returned \textsf{AMBER}: the three
\texttt{mg-365e} fixes were confirmed correct and the GREEN spine
independently re-confirmed, but a \emph{third} over-claim --- of the
identical ``a $\walsh{S}$-graded fact treated as portable'' shape as
\texttt{mg-d300}'s two findings --- survived into the revised draft.
\begin{enumerate}[label=\textup{(\Alph*)},leftmargin=2.4em]
\item \emph{The twisted-bridge null-homotopy does not port as a
  discharged input.} Construction~\ref{constr:bridge} asserted that
  the twisted-bridge identity ``ports'' to Fork~A and ``gives
  $\iota_x^{*}\simeq0$, the input the reduction uses,'' justified only
  by the observation that the $\walsh{\{x\}}$-twist is, in Fork~A, a
  cochain operation rather than a group action. That justification is
  a non sequitur: the identity needs $d$ to \emph{commute} with
  $\walsh{\{x\}}$-multiplication, a step the pre-Fork-A source
  (\texttt{mg-56b8}~\S8.1) validates \emph{only on
  $\walsh{S}$-supported cells}, so $\iota_x^{*}\simeq0$ is, at honest
  strength, a fact about the $\walsh{S}$-isotype sub-complexes
  $V_S^{*}$. Fork~A, having dropped the $(\Z/2)^n$-action
  (Remark~\ref{rem:walsh-role2}: $d$ does not preserve harmonic
  level), has no such sub-complex, so the bridge does not deliver
  $\iota_x^{*}\simeq0$ as an established Fork~A input. Establishing it
  in Fork~A is recorded here as the new sub-residual R1c
  (Remark~\ref{rem:bridge-not-ported}, Open residual~\ref{res:y1}).
  This was \texttt{mg-d300}'s ``second, compounding prong'' of its
  Finding~A, raised there in prose but omitted from its revision
  spec, so the \texttt{mg-365e} revision carried
  Construction~\ref{constr:bridge} through verbatim.
\end{enumerate}
This revision additionally carried out, per the \texttt{mg-5221}
mandate, an \emph{exhaustive sweep} of every fact in this document
ported from the pre-Fork-A Walsh / $\walsh{S}$-isotype-graded
framework, checking each for a genuine Fork~A referent. The sweep is
recorded in full in the ledger
\texttt{docs/Frankl-ForkA-Phase2-Revision2.md}; its conclusion is
that, once R1c is demoted, \emph{every} ported Walsh fact either has a
genuine Fork~A referent (the Walsh basis in its harmonic Role~2 --- a
labelling of the obstruction \emph{cochain}; see
Remark~\ref{rem:walsh-role2}), is explicitly dropped (the Walsh
grading of \emph{cohomology}; Status~\ref{status:c-block}), or is a
named residual (R1a, R1b, R1c, R3b). There is no fourth prong, and
the input list of Theorem~\ref{thm:contradiction} is complete and
final. Two minor points from \texttt{mg-5221}~\S4 are also addressed:
the base case of the simultaneous induction is now stated
(Theorem~\ref{thm:y1-reduction}, Open residual~\ref{res:y1}), and the
punctured-poset-to-$\Fstar$ extension inside
Proposition~\ref{prop:fact-one} is now explicitly attributed to the
R3 faithfulness package. As with the \texttt{mg-365e} revision, this
is an honest re-statement, not a refutation: the Fork~A strategy and
its core constructions (Theorems~\ref{thm:projection},
\ref{thm:welldefined}, \ref{thm:maschke}) are unchanged and remain
sound.
\end{status}

% =====================================================================
\section{Scope, method, and what Fork A is}
\label{sec:scope}
% =====================================================================

\subsection{The starting point}
\label{ssec:starting-point}

The cohomological arc of the union-closed program builds, for a
hypothetical minimal counterexample to Frankl, a cohomology class --- the
\emph{obstruction class} --- that two arguments force to be
simultaneously nonzero and zero. The arc was halted by its Phase-1.5
audit (work item \texttt{mg-0405}), which refuted two objects the
construction rested on:

\begin{enumerate}[label=\textup{(\alph*)},leftmargin=2.4em]
\item \textbf{the against-trace structure map.} The cohomology object
  was defined as $X^{\cap}_n=\hocolim_{(\Cn)^{\mathrm{op}}}X(\F)$, the
  homotopy colimit of a diagram over the \emph{opposite} of the trace
  category. Such a diagram requires, for each trace
  $(S,\F)\to(T,\F|_T)$, a chain map $X(\F|_T)\to X(\F)$ running
  \emph{against} the trace. Phase~1.5 exhibited an explicit object of
  $\mathcal C^{\cap}_4$ for which no such map exists (the source
  complex is a solid square, the target is four isolated points), so
  the diagram functor --- hence the object $X^{\cap}_n$ --- was
  \emph{undefined}.
\item \textbf{the $(\Z/2)^n$-action.} The Walsh grading of
  $X^{\cap}_n$, on which the vanishing and landing statements were
  phrased, required a $(\Z/2)^n$-action on the cohomology object.
  Phase~1.5 showed $(\Z/2)^n$ acts on neither the indexing category
  $\Cn$ (coordinate toggles break intersection-closure) nor the
  diagram values, so the action --- hence the Walsh decomposition ---
  was \emph{not constructed}.
\end{enumerate}

A vision-check (\texttt{mg-6edc}) read both refutations as consequences
of three documented \emph{drifts} of the build from the program's
stated goal --- ``a theorem on the cohomology of the category of
intersection-closed families, like the cohomology of the category of
posets.'' A bounded $n=3$ probe (\texttt{mg-565a}) then settled the
foundational fork:

\begin{itemize}[leftmargin=1.5em]
\item \textbf{Fork~D} (re-found on the category's \emph{bare nerve}, no
  coefficient diagram) is \emph{eliminated}: the probe computed the
  order complex of the Moore-family lattice on $[3]$ and found it
  rationally acyclic --- every reduced Betti number zero, M\"obius
  function $\mu(\hat0,\hat1)=0$ --- robustly, at $n=2$, $n=3$, and
  under both $\emptyset$-conventions. The bare nerve carries no
  cohomology and detects nothing.
\item \textbf{Fork~A} is \emph{positively endorsed}: the probe located
  the family-sensitive content precisely, in the per-family complexes
  $X(\F)$ --- a non-constant coefficient diagram over the category.
  Integrating that diagram is exactly the covariant homotopy colimit
  $\hocolim_{\Cn}X$. The probe thereby vindicated the program's
  coefficient diagram as a \emph{forced} move (the bare nerve genuinely
  has nothing) and named Fork~A as the live foundation.
\end{itemize}

\subsection{Fork A, as a precise prescription}
\label{ssec:fork-a-prescription}

Fork~A is the following four-part prescription, assembled from
\texttt{mg-0405}~\S6, the vision-check \texttt{mg-6edc}~\S\S6--7, and
the Fork-D probe \texttt{mg-565a}~\S\S6--7.

\begin{enumerate}[label=\textup{(A\arabic*)},leftmargin=2.8em]
\item \textbf{Keep} the trace category $\Cn$ and the per-family
  cubical-defect complexes $X(\F)$. The probe shows the content lives
  in the coefficient diagram $\F\mapsto X(\F)$; this is retained.
\item \textbf{Fix the variance.} Replace the (undefined) contravariant
  hocolim with the \emph{covariant} homotopy colimit
  $\hocolim_{\Cn}X$, whose structure map is the canonical trace
  \emph{projection} $\pi_T\colon X(\F)\to X(\F|_T)$. The projection
  exists and is functorial (Theorem~\ref{thm:projection}, ported from
  \texttt{mg-0405} Proposition~3.6); it is the structure map the
  refutation (a) showed the build had backwards.
\item \textbf{Drop the $(\Z/2)^n$-action} on the cohomology object
  entirely. Grade by $\sgn_{\Sym_n}$, the genuine symmetry of the
  category. Keep the Walsh apparatus only in its harmonic role --- as
  the Fourier description of the obstruction \emph{cochain} --- never
  as a group action to decompose the cohomology by.
\item \textbf{Re-establish} the obstruction class and the two-part
  contradiction on the re-founded object.
\end{enumerate}

\subsection{What Phase 1 owes, and the stance}
\label{ssec:owed}

The endorsement of Fork~A came with two named debts the probe
explicitly did \emph{not} discharge.

\begin{debt}[The re-grading]
\label{debt:regrading}
Fork~A drops the $(\Z/2)^n$-action and grades by $\sgn_{\Sym_n}$. The
vanishing and landing statements $Y_1$, $Y_2$ were phrased in Walsh
isotypes; they must be \emph{re-derived} in the $\Sym_n$-graded
language, and one must check that the closing contradiction still has
the separation it needs. This debt is treated in \S\ref{ssec:debt1}:
its structural part is discharged, and its remaining open part is
localized.
\end{debt}

\begin{debt}[Faithfulness --- the $V_3$ / wrapper-risk question]
\label{debt:faithful}
The program's stated vision has three clauses; the third ($V_3$) is
that ``the minimal counterexample differs from the generic families by
a factor relating to this cohomology'' --- i.e.\ the cohomology
\emph{genuinely detects} the counterexample. The standing risk
(``wrapper-risk'') is that the obstruction class is family-blind, in
which case the cohomology detects nothing and the contradiction is
hollow. Fork~A does not resolve this; Phase~1's task is to localize it
precisely. This debt is localized in \S\ref{ssec:debt2}.
\end{debt}

The stance throughout is default-skeptical, in the discipline the
\texttt{mg-56b8} / \texttt{mg-0405} briefs set, and with the
over-claim history of the ledger (the invalid $Y_1$ proof, the false
``$R2$'') explicitly in view: where a step closes, it is stated as a
theorem with proof; where it does not, it is isolated as a precisely
named residual, not hedged and not papered over. In particular, this
document does \emph{not} claim a clean proof of the landing statement
$Y_2$: an earlier draft of this re-foundation asserted such a proof,
and \S\ref{ssec:y2a} explains precisely why that assertion does not
hold and what is true instead.

% =====================================================================
\section{The category and the per-family complex}
\label{sec:category}
% =====================================================================

We recall the two objects Fork~A retains unchanged --- the trace
category and the per-family cubical-defect complex --- fixing notation.
Throughout, $n\ge3$, $[n]=\{1,\dots,n\}$, and all (co)homology is over
$\Q$.

\subsection{The trace category}
\label{ssec:trace-category}

\begin{definition}
\label{def:category}
The \emph{trace category} $\Cn$ has as objects the pointed
intersection-closed families: pairs $(S,\F)$ with $S\subseteq[n]$,
$\F\subseteq2^S$ closed under intersection, $\bigcup\F=S$, and
$S\in\F$. A morphism $(S,\F)\to(T,\G)$ exists precisely when
$T\subseteq S$ and $\G=\F|_T:=\{A\cap T:A\in\F\}$, the \emph{trace} of
$\F$ to $T$; this morphism is written $\tr_T$. Trace preserves
intersection-closure, composition of traces is trace
($\tr_{T'}\circ\tr_T=\tr_{T'}$ for $T'\subseteq T\subseteq S$), and
$\tr_S=\mathrm{id}$, so $\Cn$ is a genuine (finite, thin) category.
\end{definition}

This is unchanged from the pre-Fork-A development. The morphism class
remains \emph{trace}, not inclusion: Fork~D would have undone that
choice, but Fork~D is eliminated, and the Fork-D probe showed the
content lives in the coefficient diagram over $\Cn$, not in $\Cn$
itself. So $\Cn$ stays exactly as built.

\begin{notation}
\label{not:sn}
The symmetric group $\Sym_n$ acts on $[n]$ by relabelling. For
$\pi\in\Sym_n$ and $(S,\F)\in\Cn$ put $\pi\cdot(S,\F)=(\pi
S,\pi\F)$ with $\pi\F=\{\pi A:A\in\F\}$. Relabelling carries
intersection-closed families to intersection-closed families and
commutes with trace ($\pi(\F|_T)=(\pi\F)|_{\pi T}$), so $\Sym_n$ acts
on $\Cn$ by automorphisms. \emph{This action is genuine and is the
only ambient symmetry Fork~A uses.} The coordinate toggles
$\sigma_x\colon A\mapsto A\bigtriangleup\{x\}$ do \emph{not} act on
$\Cn$ --- they do not preserve intersection-closure --- which is why
Fork~A discards the $(\Z/2)^n$-action; see \S\ref{ssec:drop-z2}.
\end{notation}

\subsection{The per-family cubical-defect complex}
\label{ssec:per-family}

\begin{definition}
\label{def:xf}
For a pointed family $(S,\F)$, the \emph{cubical-defect complex}
$X(\F)$ is the cubical chain complex whose $k$-cells are the pairs
$C(A,T')$ with $A\in\F$, $T'\subseteq S\setminus A$, $|T'|=k$, and
$A\cup T''\in\F$ for every $T''\subseteq T'$ --- equivalently, the
$k$-dimensional subcubes of the cube $Q_S$ with bottom vertex $A$ and
varying directions $T'$, all $2^{k}$ of whose vertices lie in $\F$.
The differential is the standard cubical boundary $\partial$;
$C^{*}(X(\F);\Q)=\Hom(X(\F),\Q)$ is the cubical cochain complex with
coboundary $d$.
\end{definition}

$X(\F)$ records where the local cube structure ``survives'' inside
$\F$. It is a finite complex, functorial in $\F$ along traces via the
projection of the next section, and $\Sym_n$-equivariant:
$\pi\cdot X(\F)=X(\pi\F)$, with $\pi$ relabelling cells
$C(A,T')\mapsto C(\pi A,\pi T')$. This is also unchanged from the
pre-Fork-A development.

% =====================================================================
\section{The structure map: the trace projection}
\label{sec:projection}
% =====================================================================

This section establishes the single fact that makes Fork~A's
cohomology object well-defined: the trace induces a canonical chain map
\emph{in the projection direction}, making $X$ a covariant functor.
The against-trace map the pre-Fork-A build needed does not exist
(refutation (a) of \S\ref{ssec:starting-point}); the projection does.

\begin{theorem}[The trace projection is a covariant functor]
\label{thm:projection}
For every trace $\tr_T\colon(S,\F)\to(T,\F|_T)$ of $\Cn$, the
assignment
\[
  \pi_T\bigl(C(A,T')\bigr)\;=\;
  \begin{cases}
    C(A\cap T,\ T'), & T'\subseteq T,\\[2pt]
    0, & T'\not\subseteq T,
  \end{cases}
\]
defines a chain map $\pi_T\colon X(\F)\to X(\F|_T)$. It is functorial
--- $\pi_{T'}\circ\pi_T=\pi_{T'}$ for $T'\subseteq T\subseteq S$ and
$\pi_S=\mathrm{id}$ --- and $\Sym_n$-equivariant. Hence
\[
  X\colon\Cn\longrightarrow\Ch_\Q,\qquad
  (S,\F)\mapsto X(\F),\quad \tr_T\mapsto\pi_T,
\]
is a genuine \emph{covariant} functor from the trace category to
$\Q$-chain complexes.
\end{theorem}

\begin{proof}
This is the construction pinned in \texttt{mg-0405} Proposition~3.6; we
give the proof in full because it is the load-bearing positive content
of Fork~A.

\emph{Well-defined on cells.} Let $C(A,T')$ be a $k$-cell of $X(\F)$
with $T'\subseteq T$. Then $A\cap T\in\F|_T$ by definition of the
trace. Since $T'\subseteq S\setminus A$ and $T'\subseteq T$, we have
$T'\subseteq T\setminus(A\cap T)$. For every $T''\subseteq T'$,
\[
  (A\cap T)\cup T''\;=\;(A\cup T'')\cap T\;\in\;\F|_T,
\]
because $A\cup T''\in\F$ (as $C(A,T')$ is a cell of $X(\F)$) and the
trace of an element of $\F$ lies in $\F|_T$; the first equality holds
because $T''\subseteq T'\subseteq T$. So $C(A\cap T,T')$ is a genuine
$k$-cell of $X(\F|_T)$, and $\pi_T$ lands in $X(\F|_T)$.

\emph{Chain map.} $\pi_T$ is the cubical map induced by the coordinate
projection $Q_S\to Q_T$, $A\mapsto A\cap T$, on cubical chains, with
the standard normalised convention that a cell with a varying
direction in $S\setminus T$ (i.e.\ $T'\not\subseteq T$) is degenerate
and maps to $0$. A coordinate projection of cubes commutes with the
face maps: a face of $C(A,T')$ is $C(A,T'\setminus\{t\})$ or
$C(A\cup\{t\},T'\setminus\{t\})$ for $t\in T'$; if $t\in T$ these
project to the corresponding faces of $C(A\cap T,T')$, and if $t\notin
T$ the two faces in direction $t$ project to equal cells with opposite
incidence sign and cancel. Hence $\pi_T\circ\partial=\partial\circ\pi_T$.

\emph{Functorial.} For $T'\subseteq T\subseteq S$ and a cell
$C(A,U)$ with $U\subseteq T'$ one has $(A\cap T)\cap T'=A\cap T'$, and
the degenerate-to-zero convention is compatible with composition (a
direction in $S\setminus T'$ is collapsed by whichever projection
reaches it first). So $\pi_{T'}\circ\pi_T=\pi_{T'}$; and $\pi_S$ is the
identity since $A\cap S=A$.

\emph{Equivariant.} $\pi_T$ is defined from the coordinate projection
$A\mapsto A\cap T$ and the relation $T'\subseteq T$, both of which a
relabelling $\pi\in\Sym_n$ carries to the corresponding data on
$\pi T$; hence $\pi\circ\pi_T=\pi_{\pi T}\circ\pi$.
\end{proof}

\begin{remark}[Why this is the whole variance fix]
\label{rem:variance-fix}
Phase-1.5 refutation (a) was that the diagram $X$ has no functorial
structure map $X(\F|_T)\to X(\F)$ \emph{against} the trace. The
content of Theorem~\ref{thm:projection} is that it has a perfectly good
functorial structure map $\pi_T\colon X(\F)\to X(\F|_T)$ \emph{along}
the trace. The pre-Fork-A build chose the wrong variance: it wanted
$X$ to be a diagram over $(\Cn)^{\mathrm{op}}$ and asked for a map that
does not exist. Fork~A takes $X$ to be a diagram over $\Cn$ and uses
the map that does. The Phase-1.5 counterexample (the
$\mathcal C^{\cap}_4$ object whose $X(\F)$ is four isolated points
while $X(\F|_T)$ is a solid square) is no obstruction here: the
projection $\pi_T$ sends the solid square \emph{onto} (a quotient of)
the four points, which is exactly what a projection of cubes does, and
is a chain map. The minimal-lift recipe $A\mapsto\bigcap\{B\in\F:B\cap
T=A\}$ that Phase-1.5 refuted was an attempt to build the
against-trace map; Fork~A does not need it and does not use it.
\end{remark}

% =====================================================================
\section{The Fork A cohomology object}
\label{sec:object}
% =====================================================================

With $X$ a genuine covariant functor (Theorem~\ref{thm:projection}),
its homotopy colimit is a genuinely defined object. This section
constructs it and its cohomology, and records that the construction
clears Phase-1.5 refutation (a).

\subsection{The Bousfield--Kan double complex}
\label{ssec:bk}

\begin{definition}
\label{def:bk}
The \emph{Fork~A cohomology object} is the homotopy colimit
$\Xbb_n:=\hocolim_{\Cn}X$ of the covariant functor $X$ of
Theorem~\ref{thm:projection}. For cohomology over $\Q$ we work with the
dual cochain Bousfield--Kan double complex
\[
  \BK^{p,q}\;=\;\bigoplus_{\F_0\xrightarrow{\ \tr\ }\F_1\xrightarrow{\ \tr\ }
  \cdots\xrightarrow{\ \tr\ }\F_p\ \text{in }\Cn}
  C^{q}\bigl(X(\F_0);\Q\bigr),
  \qquad p,q\ge0,
\]
indexed by the bar strings of length $p$ in $\Cn$, with the cochain
value taken at the \emph{first} object $\F_0$ of each string (the sum
is finite: $\Cn$ has finitely many objects and morphisms). It carries
two differentials:
\begin{itemize}[leftmargin=1.6em]
\item the \emph{vertical} differential $d\colon\BK^{p,q}\to\BK^{p,q+1}$,
  the cubical coboundary applied to each $C^{q}(X(\F_0))$;
\item the \emph{horizontal} (bar) differential
  $\delta_{\BK}\colon\BK^{p,q}\to\BK^{p+1,q}$,
  $\delta_{\BK}=\sum_{j=0}^{p+1}(-1)^{j}d_j^{*}$, where for a string
  $\F_0\to\cdots\to\F_{p+1}$:
  \begin{itemize}[leftmargin=1.4em]
  \item $d_0^{*}$ drops $\F_0$ and pulls the value at $\F_1$ back to
    $\F_0$ along the projection of the first trace, via
    $\pi_{T_1}^{*}\colon C^{q}(X(\F_1))\to C^{q}(X(\F_0))$;
  \item $d_j^{*}$ for $1\le j\le p$ composes the $j$-th and
    $(j{+}1)$-st traces (no change of cochain value);
  \item $d_{p+1}^{*}$ drops the last object $\F_{p+1}$ (no change of
    cochain value, which sits at $\F_0$).
  \end{itemize}
\end{itemize}
The \emph{Fork~A cohomology} is
$H^{*}(\Cn;X):=H^{*}\bigl(\Tot\BK,\,D\bigr)$, the cohomology of the
total complex with total differential $D=d+(-1)^{q}\delta_{\BK}$.
\end{definition}

\begin{theorem}[The Fork A object is genuinely defined]
\label{thm:welldefined}
$\Tot\BK$ with the differential $D$ of Definition~\ref{def:bk} is a
genuine cochain complex: $D^2=0$. Consequently $H^{*}(\Cn;X)$ is a
well-defined graded $\Q$-vector space.
\end{theorem}

\begin{proof}
$d^2=0$ is the cubical cochain identity. $\delta_{\BK}^{2}=0$ is the
cosimplicial (bar) identity: it holds for the cochain bar construction
of \emph{any} covariant functor into an abelian category, and $X$ is
such a functor by Theorem~\ref{thm:projection}. The only component of
$\delta_{\BK}$ that involves the diagram values --- the face $d_0^{*}$
--- uses $\pi_{T}^{*}$, the cochain dual of the genuine chain map
$\pi_T$; the bar identity $\delta_{\BK}^2=0$ then reduces, at the
$d_0^{*}d_0^{*}$ term, to the functoriality
$\pi_{T'}\circ\pi_T=\pi_{T'}$, which Theorem~\ref{thm:projection}
supplies. Finally $d$ and $\delta_{\BK}$ commute up to the standard
sign because each $\pi_T^{*}$ is a cochain map ($\pi_T$ is a chain map,
Theorem~\ref{thm:projection}), so $d\,\pi_T^{*}=\pi_T^{*}\,d$. Hence
$D^2=d^2\pm(d\,\delta_{\BK}-\delta_{\BK}\,d)+\delta_{\BK}^2=0$.
\end{proof}

\begin{status}
\label{status:b-block}
Theorem~\ref{thm:welldefined} \emph{dodges Phase-1.5 refutation (a)
by construction.} The pre-Fork-A object $X^{\cap}_n$ was undefined
because its diagram functor required a structure map that does not
exist; the Fork~A object $H^{*}(\Cn;X)$ is defined because its
diagram functor requires only $\pi_T^{*}$, which exists for free as
the dual of the projection. No new mathematics is needed for this:
Theorem~\ref{thm:projection} is \texttt{mg-0405}'s own
Proposition~3.6, and \texttt{mg-0405}~\S6 already named ``adopt the
covariant hocolim'' as Fork~A. Phase~1 carries it out and verifies
$D^2=0$ explicitly. This step is \textsf{GREEN}.
\end{status}

\subsection{What this object is}
\label{ssec:what-it-is}

$H^{*}(\Cn;X)$ is the cohomology of the trace category $\Cn$ with
coefficients in the covariant cubical-defect coefficient system
$X(-)$. It is \emph{not} the cohomology of the bare nerve of $\Cn$;
the Fork-D probe showed the bare nerve is rationally acyclic and
detects nothing, so a non-constant coefficient system is forced, and
$X(-)$ is that system. It is \emph{not} the cohomology of a diagram
over $(\Cn)^{\mathrm{op}}$; that was the refuted variance. It is the
homotopy colimit of the genuine covariant diagram --- the object the
Fork-D probe endorsed in so many words: ``Integrating that diagram
over the category is exactly $\hocolim_{\Cn}X(\F)$, and its cohomology
$H^{*}(\Cn;X(-))$ is the cohomology of the category with coefficients''
(\texttt{mg-565a}~\S6.2).

\begin{remark}[A faithfulness flag, recorded now, discharged in \S\ref{ssec:debt2}]
\label{rem:hocolim-flag}
The obstruction class of \S\ref{sec:obstruction} is a
\emph{gluing}-type invariant --- it measures whether local
rare-coordinate data extend to a global section. Gluing is a limit
phenomenon, and a referee may reasonably ask whether such a class is
correctly housed in the cohomology of a homotopy \emph{colimit}. The
program's prescription --- vision-check, Fork-D probe, and
\texttt{mg-0405}~\S6 alike --- names the covariant hocolim, and Phase~1
follows it; but the hocolim-versus-holim question is genuine and is
folded, honestly, into the faithfulness debt (Owed
debt~\ref{debt:faithful}, \S\ref{ssec:debt2}). It is flagged here so
that it is not mistaken for settled.
\end{remark}

% =====================================================================
\section{The re-grading: from $(\Z/2)^n$ to $\sgn_{\Sym_n}$}
\label{sec:regrading}
% =====================================================================

This section discharges the structural half of Owed
debt~\ref{debt:regrading}: it drops the $(\Z/2)^n$-action,
re-establishes the grading by the genuine $\Sym_n$-action, and
identifies the sphere class. The statement-level half --- re-deriving
$Y_1$, $Y_2$ --- is \S\ref{sec:contradiction}.

\subsection{Dropping the $(\Z/2)^n$-action}
\label{ssec:drop-z2}

\begin{status}
\label{status:c-block}
Fork~A \emph{drops} the $(\Z/2)^n$-action on the cohomology object.
This is not a repair of Phase-1.5 refutation (b) --- it is its
\emph{dissolution}. Refutation (b) said the $(\Z/2)^n$-action does not
assemble: $(\Z/2)^n$ acts on neither $\Cn$ nor the diagram values, so
the Walsh decomposition of the cohomology was never constructed.
Fork~A's response is not to construct it but to observe it is not
needed for the grading. The grading the program actually wants --- the
sphere carries $\sgn$ --- is a statement about the $\Sym_n$-action,
which is genuine (Notation~\ref{not:sn}). With no $(\Z/2)^n$-action in
the construction, \emph{there is nothing left for refutation (b) to
refute.} This step is \textsf{GREEN}. \emph{Caveat, recorded honestly:}
dropping the $(\Z/2)^n$-action has a downstream cost on \emph{both}
halves of the contradiction. The pre-Fork-A landing argument $Y_2$
used the $(\Z/2)^n$-action (it split the obstruction off by Walsh
isotype); dropping the action means $Y_2$ must be genuinely
re-derived, not merely re-graded (see~\S\ref{ssec:y2a}, residual~R2).
And the pre-Fork-A $Y_1$ cofiber-LES reduction split
isotype-by-isotype using the same $(\Z/2)^n$-action, which --- being
abelian, with a quotient acting on the smaller cube --- descends to
that cube; $\Sym_n$ has no such quotient, so the $Y_1$ reduction
acquires an $\Sym_{n-1}\to\Sym_n$ branching step
(see~\S\ref{ssec:y1a}, Remark~\ref{rem:y1-branching}, residual~R1b).
Dropping $(\Z/2)^n$ \emph{also} removes the $\walsh{S}$-isotype
sub-complexes $V_S^{*}$ on which the pre-Fork-A twisted-bridge
null-homotopy $\iota_x^{*}\simeq0$ lived, so that null-homotopy
ceases to be a discharged input and becomes the sub-residual~R1c
(see~\S\ref{ssec:y1a}, Remark~\ref{rem:bridge-not-ported}). Phase~1
records all of these costs.
\end{status}

The Walsh apparatus is not discarded wholesale --- only its
\emph{Role~1} use, as a group action to split the cohomology by. Its
\emph{Role~2} use survives, with one honest qualification:

\begin{remark}[Walsh, Role 2: the harmonic description and its limits]
\label{rem:walsh-role2}
For each ground set $S$, the Walsh characters
$\walsh{U}\colon2^{S}\to\{\pm1\}$, $\walsh{U}(A)=(-1)^{|U\cap A|}$,
form an orthonormal basis of $\Q$-valued functions on the cube $2^S$.
This is pure harmonic analysis: it requires no group \emph{action} on
any complex, only the function basis. In particular the bias is, up to
scale and sign, the level-$1$ Walsh--Fourier coefficient of the family
indicator, $\bias_x(\F)=-2^{n}\widehat{\mathbf1_\F}(\{x\})$, and the
obstruction cochain (\S\ref{sec:obstruction}) is built from exactly
this level-$1$ data. Fork~A keeps the Walsh basis in precisely this
role --- the harmonic description of the obstruction \emph{cochain}.
\emph{The honest qualification:} the \emph{harmonic level} of a
cochain is well-defined, but it is \emph{not} preserved by the cubical
coboundary $d$ on the defect complex $X(\F)$ --- Phase~1.5 established
exactly this (the single-family Walsh operators are not cochain maps).
So Walsh level grades \emph{cochains}, not \emph{cohomology}. This is
not a minor point: it is the reason the landing statement $Y_2$ cannot
be read off the harmonic level of the obstruction cochain, and it is
the crux of \S\ref{ssec:y2a}.
\end{remark}

\subsection{The $\Sym_n$-grading by Maschke}
\label{ssec:maschke}

\begin{theorem}[The $\Sym_n$-isotype decomposition]
\label{thm:maschke}
The action of $\Sym_n$ on $\Cn$ (Notation~\ref{not:sn}) and its
equivariant action on the diagram $X$ (Theorem~\ref{thm:projection})
induce an action of $\Sym_n$ on $\Tot\BK$ by cochain automorphisms.
Hence each $H^{k}(\Cn;X)$ is a finite-dimensional $\Q[\Sym_n]$-module,
and by Maschke's theorem it decomposes canonically into isotypic
components
\[
  H^{k}(\Cn;X)\;=\;\bigoplus_{\lambda\,\vdash\,n}
  H^{k}(\Cn;X)_{[\lambda]},
\]
one summand for each irreducible $\Q$-representation $\lambda$ of
$\Sym_n$. Moreover the total complex itself splits, $\Tot\BK=
\bigoplus_{\lambda}(\Tot\BK)_{[\lambda]}$, as a direct sum of
sub-cochain-complexes, and $H^{k}(\Cn;X)_{[\lambda]}=
H^{k}\bigl((\Tot\BK)_{[\lambda]}\bigr)$.
\end{theorem}

\begin{proof}
$\Sym_n$ acts on the index set of $\BK^{p,q}$ --- the bar strings of
$\Cn$ --- by relabelling, and on each summand $C^{q}(X(\F_0))$ by the
equivariance $\pi\cdot X(\F_0)=X(\pi\F_0)$. The two differentials are
$\Sym_n$-equivariant: $d$ because the cubical coboundary is natural
under relabelling, and $\delta_{\BK}$ because $\pi_T$ is equivariant
(Theorem~\ref{thm:projection}) and the bar faces are defined from the
category structure, which $\Sym_n$ preserves. So $\Sym_n$ acts on
$\Tot\BK$ by cochain automorphisms. $\Sym_n$ is finite and $\Q$ has
characteristic $0$, so Maschke applies: $\Q[\Sym_n]$ is semisimple,
every $\Q[\Sym_n]$-module splits canonically into isotypes, and an
equivariant map (here $D$) is block-diagonal for that splitting.
Splitting $\Tot\BK$ degreewise into isotypes therefore splits it as a
complex, and cohomology commutes with the finite direct sum.
\end{proof}

\begin{remark}[This is the genuine analogue of the old Walsh decomposition]
\label{rem:maschke-analogue}
The pre-Fork-A development had a ``Walsh decomposition theorem'' whose
content was Maschke applied to a $(\Z/2)^n$-action.
Theorem~\ref{thm:maschke} is Maschke applied to the $\Sym_n$-action
--- the decisive difference being that the $\Sym_n$-action is genuine
(Notation~\ref{not:sn}) whereas the $(\Z/2)^n$-action was not
(Phase-1.5 refutation (b)). The old decomposition was Maschke for a
group that does not act; Theorem~\ref{thm:maschke} is Maschke for the
group that does. This is the precise sense in which the re-grading is
sound. It does \emph{not} follow that every statement phrased in Walsh
isotypes translates to an $\Sym_n$-isotype statement: the two gradings
are different decompositions of the cohomology (when both are defined),
not nested. Re-deriving $Y_1$ and $Y_2$ in the $\Sym_n$ grading is
genuine work, carried out --- and honestly assessed --- in
\S\ref{sec:contradiction}.
\end{remark}

\subsection{The sphere class}
\label{ssec:sphere-class}

The program's structural target --- ``it has spherical cohomology,''
the vision's clause $V_2$ --- is, on the cube side, that
$H^{*}(\Cn;X)$ is concentrated in the top degree $n-1$ and
one-dimensional there. Under the $\Sym_n$-grading this top class is
forced to be the sign:

\begin{proposition}[The sphere class is the sign isotype]
\label{prop:sphere-sgn}
Suppose $H^{*}(\Cn;X)$ is rationally a homotopy $(n-1)$-sphere:
$\widetilde H^{k}(\Cn;X)=0$ for $0\le k\le n-2$ and $\widetilde
H^{n-1}(\Cn;X)$ one-dimensional. Suppose moreover that the top class
is realised $\Sym_n$-equivariantly as the orientation class of an
$(n-1)$-sphere on which $\Sym_n$ acts by permuting $n$ coordinate
directions --- equivalently, that some transposition acts on
$\widetilde H^{n-1}(\Cn;X)$ by $-1$. Then, as a $\Q[\Sym_n]$-module,
\[
  \widetilde H^{n-1}(\Cn;X)\;\cong\;\sgn_{\Sym_n}.
\]
\end{proposition}

\begin{proof}
A one-dimensional $\Q$-representation of $\Sym_n$ is either the trivial
or the sign representation, and the two are distinguished exactly by
whether a transposition acts by $+1$ or by $-1$. The added hypothesis
fixes the latter: a transposition, permuting two of the $n$ coordinate
directions of the equivariant sphere, acts on the orientation class by
the degree of the induced self-map, which is the sign of that
transposition, namely $-1$. Hence the one-dimensional module is
$\sgn_{\Sym_n}$, not the trivial one. (This is the cube-side analogue
of the F-series fact that $\widetilde
H^{n-2}(\Delta(\mathrm{PPF}_n);\Q)\cong\sgn_{\Sym_n}$: the top
cohomology of an $\Sym_n$-symmetric sphere is the orientation
character, the sign.)
\end{proof}

\begin{remark}[Why the strengthened hypothesis]
\label{rem:sphere-sgn-hyp}
The bare dimension count --- $\widetilde H^{n-1}(\Cn;X)$
one-dimensional --- does \emph{not} by itself force the sign: a
one-dimensional $\Q[\Sym_n]$-module may a priori be trivial, and
ruling that out requires the genuine $\Sym_n$-action on the class, not
an analogy with coordinate spheres. An earlier draft's proof appealed
to ``a sphere built $\Sym_n$-symmetrically from $n$ coordinate
directions'' without that structure being part of the hypothesis;
Proposition~\ref{prop:sphere-sgn} now takes the equivariant
coordinate-sphere structure (equivalently, the action of one
transposition) as an explicit hypothesis. It is not load-bearing for
the contradiction: Theorem~\ref{thm:contradiction} does not use
Proposition~\ref{prop:sphere-sgn} but input~(v)/$Y_1$-A
(Conjecture~\ref{conj:y1a}), which asserts $\widetilde
H^{n-1}(\Cn;X)\cong\sgn_{\Sym_n}$ directly. Proposition~%
\ref{prop:sphere-sgn} records only \emph{what} the top class is, and
on \emph{what} hypothesis; closing $Y_1$-A is what would supply that
hypothesis genuinely.
\end{remark}

\begin{remark}[The $\sgn$ was always an $\Sym_n$-fact]
\label{rem:sgn-always-sn}
The pre-Fork-A target was $\widetilde H^{n-1}\cong\walsh{[n]}\boxtimes
\sgn_{\Sym_n}$. The $\walsh{[n]}$ tensor factor was the
$(\Z/2)^n$-decoration; the $\sgn_{\Sym_n}$ factor was always carried by
the $\Sym_n$-action on the one-dimensional multiplicity space.
Proposition~\ref{prop:sphere-sgn} recovers exactly the $\sgn_{\Sym_n}$
factor and discards exactly the $\walsh{[n]}$ factor. So dropping
$(\Z/2)^n$ costs nothing at the sphere class: the sign was never a
$(\Z/2)^n$-datum. This is the first concrete confirmation that the
re-grading (Owed debt~\ref{debt:regrading}) is sound at the structural
level.
\end{remark}

We do \emph{not} claim the sphere-concentration hypothesis of
Proposition~\ref{prop:sphere-sgn}; that is the $Y_1$ statement, and it
is open (\S\ref{ssec:y1a}). Proposition~\ref{prop:sphere-sgn} says only
\emph{what} the top class is, \emph{if} the cohomology is spherical.

% =====================================================================
\section{The obstruction class, re-founded}
\label{sec:obstruction}
% =====================================================================

This section re-establishes the obstruction class on the Fork~A object.
The construction is verbatim from the pre-Fork-A development --- the
\v{C}ech double complex of the obstruction sheaf --- because nothing in
that construction used the refuted variance or the refuted
$(\Z/2)^n$-action; it used only the trace poset below a minimal
counterexample and the bias data. What changes is only the object the
class is promoted into: $H^{n-1}(\Cn;X)$ in place of the undefined
$X^{\cap}_n$.

\subsection{The minimal counterexample and the trace cover}
\label{ssec:trace-cover}

Work in the intersection-closed dialect, in which Frankl reads: every
intersection-closed $\F\neq\{[n]\}$ has a \emph{rare} coordinate, some
$x$ with $\bias_x(\F)\le0$. Suppose Frankl fails.

\begin{definition}
\label{def:minimal}
A \emph{minimal counterexample} $\Fstar$ is an intersection-closed
family on $[n]$ with $\bias_x(\Fstar)>0$ for every $x$ (no rare
coordinate), minimising $|\Fstar|$, then $n$.
\end{definition}

\begin{proposition}[Minimality--element]
\label{prop:minimality}
Every proper trace $\Fstar|_T$, $T\subsetneq[n]$, has a rare
coordinate: some $x\in T$ with $\bias_x(\Fstar|_T)\le0$.
\end{proposition}

\begin{proof}
$\Fstar|_T$ is a strictly smaller intersection-closed family; if it
had no rare coordinate it would be a smaller counterexample,
contradicting the minimality of $\Fstar$.
\end{proof}

Proposition~\ref{prop:minimality} furnishes a partial rare-coordinate
assignment, defined on every proper trace of $\Fstar$ but \emph{not}
on $\Fstar$ itself. The obstruction measures the failure to extend it.
For $x\in[n]$ let
\[
  U_x\;=\;\bigl\{(T,\G)\neq\Fstar\ :\ x\in T,\ \bias_x(\G)\le0\bigr\}
\]
be the stratum of proper traces in which $x$ is rare. By
Proposition~\ref{prop:minimality} the strata $\{U_x\}_{x\in[n]}$ cover
the punctured trace poset $\Cn(\Fstar)\setminus\{\Fstar\}$, where
$\Cn(\Fstar)$ is the full subcategory of $\Cn$ on the traces of
$\Fstar$.

\subsection{The obstruction double complex}
\label{ssec:obs-complex}

\begin{definition}
\label{def:obs-sheaf}
On the stratum $U_x$ place the \emph{local bias cochain}
$\widehat b_x=b_x\cdot[\walsh{\{x\}}]$, where $b_x(\G)=-\bias_x(\G)\ge0$
is the (non-negative) defect and $[\walsh{\{x\}}]$ records that the
datum is carried by the level-$1$ Walsh character $\walsh{\{x\}}$
(Remark~\ref{rem:walsh-role2}, Role~2 --- this is a labelling of the
harmonic type of the cochain, not a group action). On a pairwise
overlap $U_x\cap U_y$ place the \emph{mismatch cochain}
$m_{xy}=\widehat b_x-\widehat b_y$. These data form the
\emph{obstruction sheaf} $\fobs$ on $\Cn(\Fstar)\setminus\{\Fstar\}$
and assemble into the \v{C}ech double complex of the cover
$\mathfrak U=\{U_x\}_{x\in[n]}$,
\[
  \check C^{p,q}(\mathfrak U;\fobs)\;=\;
  \bigoplus_{x_0<x_1<\cdots<x_p}\fobs^{q}\bigl(U_{x_0}\cap\cdots\cap
  U_{x_p}\bigr),
\]
with \v{C}ech differential $\check\delta$ (alternating sum over
omitted indices) and internal cochain differential $d$.
\end{definition}

\begin{lemma}[Cocycle]
\label{lem:cocycle}
On every triple overlap $m_{xy}+m_{yz}+m_{zx}=0$, so the family
$\{m_{xy}\}$ is a \v{C}ech $1$-cocycle. It is the \v{C}ech coboundary
$\check\delta^{0}\{\widehat b_x\}$ of the $0$-cochain
$\{\widehat b_x\}$.
\end{lemma}

\begin{proof}
$m_{xy}+m_{yz}+m_{zx}=(\widehat b_x-\widehat b_y)+(\widehat
b_y-\widehat b_z)+(\widehat b_z-\widehat b_x)=0$, and by definition
$(\check\delta^{0}\{\widehat b_x\})_{xy}=\widehat b_x-\widehat
b_y=m_{xy}$.
\end{proof}

The $\{\widehat b_x\}$ are genuine local sections of $\fobs$ but do
\emph{not} agree on overlaps --- that is the whole point: a minimal
counterexample is precisely a family for which they cannot be made to
agree. So $\{\widehat b_x\}$ is not a global section, and the mismatch
$\{m_{xy}\}$, though a $\check\delta$-cocycle, need not be a coboundary
\emph{of a global section}. The obstruction is exactly that
discrepancy.

\begin{definition}[The obstruction class]
\label{def:ob}
The mismatch cocycle $\{m_{xy}\}$, promoted through the edge map of the
\v{C}ech-to-cohomology spectral sequence of the double complex
$\check C^{*,*}(\mathfrak U;\fobs)$, defines the \emph{obstruction
class}
\[
  \ob(\Fstar)\;\in\;\widetilde H^{n-1}(\Cn;X).
\]
This membership is \emph{conditional}: it presupposes that the edge map
exists and lands in degree $n-1$ (parts R3a, R3b of Owed
debt~\ref{debt:faithful-localized}). Unconditionally, Phase~1
constructs $\{m_{xy}\}$ as a \v{C}ech $1$-cocycle
(Lemma~\ref{lem:cocycle}); see Remark~\ref{rem:promotion}.
\end{definition}

\begin{remark}[The promotion, stated honestly]
\label{rem:promotion}
The spectral sequence of the double complex
$\check C^{*,*}(\mathfrak U;\fobs)$ has $E_2^{p,q}=\check
H^{p}(\mathfrak U;\mathcal H^{q}(\fobs))$ and abuts to its own total
cohomology. Promoting the mismatch class from there into
$H^{*}(\Cn;X)$ requires a further comparison map, and that comparison
crosses a genuine object mismatch: the \v{C}ech double complex lives
on the punctured trace poset $\Cn(\Fstar)\setminus\{\Fstar\}$, whereas
$H^{*}(\Cn;X)$ is the cohomology of the \emph{whole} category $\Cn$.
The mismatch cocycle sits in \v{C}ech degree $p=1$, and the total
degree $n-1$ asserted for $\ob(\Fstar)$ is the cube-side dimension
count inherited from the pre-Fork-A development --- a count performed
in the now-defunct object $X^{\cap}_n$, not re-verified in the
genuinely different Fork~A object. The precise pinning of this edge
map --- its \emph{existence}, its \emph{degree behaviour}, and
(crucially) its \emph{injectivity} --- is \emph{not} discharged here;
it is exactly the three-part Owed debt~\ref{debt:faithful-localized}
(R3a existence, R3b degree, R3c injectivity) and is localized in
\S\ref{ssec:debt2}. Consequently $\ob(\Fstar)$ is, in Phase~1, a
well-defined \emph{\v{C}ech cocycle}; its promotion into a genuine
element of $\widetilde H^{n-1}(\Cn;X)$ is conditional on R3a and R3b.
What Phase~1 does establish, and uses, is only that the \v{C}ech
double complex --- and any edge map natural in the ground set --- is
$\Sym_n$-equivariant, since every ingredient (the cover indexed by
$[n]$, the sheaf $\fobs$, the differentials) is natural in $[n]$.
\end{remark}

% =====================================================================
\section{The two-part contradiction, re-established}
\label{sec:contradiction}
% =====================================================================

The intended proof of Frankl is a collision: the obstruction class is
forced nonzero by minimality (Fact One) and zero by the structure of
the cohomology (Fact Two). This section re-establishes the
contradiction on the Fork~A object as a \emph{conditional} theorem,
and is scrupulous about which of its inputs are established and which
are open.

\subsection{Fact One: the obstruction is nonzero}
\label{ssec:fact-one}

\begin{proposition}[Fact One, conditional]
\label{prop:fact-one}
Suppose the \v{C}ech-to-cohomology edge map of Definition~\ref{def:ob}
exists and carries the mismatch class into degree $n-1$ (so that
$\ob(\Fstar)$ is a well-defined element of $\widetilde
H^{n-1}(\Cn;X)$), and suppose it is injective on the class of the
mismatch cocycle. Then $\ob(\Fstar)\neq0$ in $\widetilde
H^{n-1}(\Cn;X)$.
\end{proposition}

\begin{proof}
Suppose $\ob(\Fstar)=0$. If the edge map is injective, this forces the
mismatch cocycle $\{m_{xy}\}$ to be a \v{C}ech coboundary \emph{of a
global section}: the local data $\{\widehat b_x\}$ glue. A global
section of $\fobs$ is, strictly, a single rare-coordinate datum valid
on the \emph{punctured} poset $\Cn(\Fstar)\setminus\{\Fstar\}$ on
which $\fobs$ is defined (Definition~\ref{def:obs-sheaf}); extending
it to a datum valid at the top object $\Fstar$ \emph{itself} is the
punctured-poset-to-$\Fstar$ promotion that
Remark~\ref{rem:promotion} folds into the R3 faithfulness package
(specifically the edge-map \emph{existence} part R3a, which is what
bridges the punctured poset and the whole category). Granting the
standing hypothesis that the edge map exists, that promotion produces
a coordinate $x_\ast$ with $\bias_{x_\ast}(\Fstar)\le0$. But a
minimal counterexample has $\bias_x(\Fstar)>0$ for every $x$
(Definition~\ref{def:minimal}). Contradiction; hence
$\ob(\Fstar)\neq0$.
\end{proof}

\begin{status}
\label{status:fact-one}
Fact One is \emph{conditional on the edge map} --- on its existence,
its degree-$(n-1)$ behaviour, and its injectivity, the three parts of
residual R3 (Owed debt~\ref{debt:faithful-localized},
\S\ref{ssec:debt2}). The \emph{injectivity} is not a defect introduced
by Fork~A --- the pre-Fork-A development carried exactly the same
conditionality, recorded there as ``the load-bearing content is the
equivalence $\ob=0\iff$ a global section exists, and that equivalence
is the assertion that the obstruction theory is faithful.'' The
\emph{existence} and \emph{degree} parts, by contrast, are the honest
residue of the fact that Definition~\ref{def:ob} promotes the mismatch
cocycle through an edge map that Remark~\ref{rem:promotion} constructs
only at the level of the \v{C}ech double complex, not as a pinned map
into $H^{*}(\Cn;X)$ --- so Fact One, like Theorem~\ref{thm:contradiction},
presupposes that $\ob(\Fstar)$ is a genuine element of $\widetilde
H^{n-1}(\Cn;X)$ in the first place. In particular, the
punctured-poset-to-$\Fstar$ extension that
Proposition~\ref{prop:fact-one}'s proof performs --- passing from a
global section of $\fobs$ on $\Cn(\Fstar)\setminus\{\Fstar\}$ to a
rare-coordinate datum valid at $\Fstar$ --- is itself part of the R3a
edge-map package (Remark~\ref{rem:promotion}), not a step Phase~1
takes for free. Phase~1 discharges none of the three.
\end{status}

\subsection{Fact Two, part 1: the landing ($Y_2$-A)}
\label{ssec:y2a}

The landing statement asks where in $H^{n-1}(\Cn;X)$ the obstruction
class sits. In the pre-Fork-A development this was $Y_2$: ``$\ob$
lands in the level-$1$ Walsh isotypes,'' and its mechanism was that
the additive, $(\Z/2)^n$-equivariant differentials of the \v{C}ech
double complex preserve Walsh isotype. \textbf{That mechanism does not
survive Fork~A}: it used $(\Z/2)^n$-equivariance, and Fork~A has no
$(\Z/2)^n$-action (Status~\ref{status:c-block}). The landing must
therefore be genuinely re-derived in the $\Sym_n$ grading. We state
honestly what re-derives, and what does not.

\subsubsection*{What does re-derive: the structural mechanism}

\begin{proposition}[$Y_2$-A, structural mechanism]
\label{prop:y2a-mech}
The obstruction class $\ob(\Fstar)$ has the following two properties.
\begin{enumerate}[label=\textup{(\arabic*)},leftmargin=2.4em]
\item \textup{(it factors through mismatches)} $\ob(\Fstar)$ depends on
  the local bias data $(\widehat b_x)_{x\in[n]}$ \emph{only through the
  differences} $\widehat b_x-\widehat b_y$. Equivalently
  (Lemma~\ref{lem:cocycle}), it factors through the \v{C}ech coboundary
  map $\check\delta^{0}$, hence is insensitive to any datum that is
  already a global section of $\fobs$.
\item \textup{(it is harmonically level-$1$ at the cochain level)} The
  obstruction \emph{cochain} --- the mismatch $\{m_{xy}\}$ and every
  cochain in the \v{C}ech direction built from it by the additive
  \v{C}ech differential --- is supported on the level-$1$ Walsh
  characters $\walsh{\{x\}}$; the \v{C}ech differential is additive
  and carries out no multiplication by a Walsh character, so it does
  not raise the harmonic level of a cochain.
\end{enumerate}
Both properties are $\Sym_n$-equivariant statements and hold with no
group action on the cohomology object.
\end{proposition}

\begin{proof}
(1) By Definition~\ref{def:ob} the obstruction is the promotion of the
mismatch $\{m_{xy}\}=\{\widehat b_x-\widehat b_y\}$; every $m_{xy}$ is
a difference. By Lemma~\ref{lem:cocycle} the mismatch is
$\check\delta^{0}\{\widehat b_x\}$, and $\check\delta^{0}$ annihilates
exactly the global sections (the $0$-cochains agreeing on overlaps).
(2) Each $\widehat b_x=b_x\cdot[\walsh{\{x\}}]$ is level-$1$ harmonic
by Definition~\ref{def:obs-sheaf}; the \v{C}ech differential
$\check\delta$ is an alternating sum of restrictions, an additive
operation involving no Walsh multiplication, so it preserves the
level-$1$ harmonic support. Both constructions are natural in the
ground set $[n]$, hence $\Sym_n$-equivariant.
\end{proof}

\subsubsection*{What does not re-derive for free: the separation from the sphere}

The closing contradiction needs more than Proposition~\ref{prop:y2a-mech}.
It needs the obstruction's cohomology \emph{class} to lie in a part of
$H^{n-1}(\Cn;X)$ that the vanishing statement $Y_1$-A can kill ---
concretely, a part \emph{disjoint from the $\sgn$ sphere-class
isotype}, since the sphere class is the one piece $Y_1$-A leaves
nonzero (\S\ref{ssec:sphere-class}). In the pre-Fork-A development this
separation was free: $\ob$ sat at Walsh level $1$, the sphere class at
Walsh level $n$, and Walsh level separated them. Fork~A has dropped
Walsh level as a grading of the cohomology. The honest finding is:

\begin{openresidual}[$Y_2$-A: the separation residual --- R2]
\label{res:y2}
It is \emph{not} established by Phase~1 that $\ob(\Fstar)$ avoids the
$\sgn_{\Sym_n}$ sphere-class isotype of $H^{n-1}(\Cn;X)$. The two
properties of Proposition~\ref{prop:y2a-mech} do not deliver it:
\begin{itemize}[leftmargin=1.5em]
\item \emph{Harmonic level is not a cohomology grading.} The
  obstruction cochain is level-$1$ harmonic
  (Proposition~\ref{prop:y2a-mech}(2)), but the cubical coboundary $d$
  does not preserve harmonic level on the defect complex
  (Remark~\ref{rem:walsh-role2}); so ``the cochain is level-$1$'' does
  not imply ``the class avoids the top-harmonic / sphere part.''
\item \emph{Factoring through mismatches constrains, but does not
  pin, the $\Sym_n$-isotype.} By
  Proposition~\ref{prop:y2a-mech}(1) and Maschke
  (Theorem~\ref{thm:maschke}), $\ob(\Fstar)$ lies in the image of an
  $\Sym_n$-equivariant map out of $\im\check\delta^{0}$, so its
  $\Sym_n$-isotype content is contained in that of
  $\im\check\delta^{0}\subseteq\check C^{1,*}$. Whether
  $\im\check\delta^{0}$ contains the $\sgn_{\Sym_n}$ isotype is a
  representation-theoretic property of the obstruction sheaf $\fobs$:
  by Frobenius reciprocity for the cover indexed by $[n]$, it reduces
  to whether the coefficient $\fobs^{q}(U_x)$ carries the relevant
  sign character of the point stabiliser $\Sym_{n-1}$. The fine
  structure of $\fobs^{q}$ is not pinned in Phase~1.
\end{itemize}
\textbf{$Y_2$-A reduces to the following named residual:}
\emph{$\ob(\Fstar)$ has no $\sgn_{\Sym_n}$-component in
$H^{n-1}(\Cn;X)$} --- equivalently, the level-$1$-harmonic obstruction
class is not cohomologous to a nonzero multiple of the
top-harmonic sphere class. This is a representation-theoretic question
about the obstruction sheaf $\fobs$, and it is naturally a
\emph{faithfulness} question: if $\ob$ had a $\sgn$-component it would
be partly the family-blind orientation class, which is the precise
shape of the wrapper-risk.
\end{openresidual}

\begin{status}
\label{status:y2a}
$Y_2$-A is \textsf{AMBER --- open}. The structural mechanism re-derives
in the $\Sym_n$ grading with no group action on the cohomology
(Proposition~\ref{prop:y2a-mech}); this is the genuine portable
content of the pre-Fork-A $Y_2$. But the conclusion the contradiction
needs --- that $\ob$ avoids the sphere isotype --- is \emph{not} free,
because the pre-Fork-A $Y_2$ obtained it from the Walsh-level grading
of cohomology, which Fork~A correctly discards. Phase~1 isolates the
residual precisely (Open residual~\ref{res:y2}). An earlier draft of
this re-foundation asserted a clean theorem ``$\ob$ is
standard-isotypic, hence $\sgn$-free''; that assertion is
\emph{withdrawn}: it tacitly assumed the coefficient $\fobs^{q}$ is
$\Sym_n$-trivial, which is not justified, and the honest statement is
the residual above. Recording this correction is part of the Phase-1
deliverable: the over-claim discipline that caught the invalid $Y_1$
proof and the false ``$R2$'' applies here too.
\end{status}

\subsection{Fact Two, part 2: the vanishing ($Y_1$-A)}
\label{ssec:y1a}

The closing contradiction needs the obstruction's home to vanish.
Granting (R2) that $\ob$ avoids the sphere isotype, the home is the
non-$\sgn$ part of $\widetilde H^{n-1}(\Cn;X)$. So the Fork~A
vanishing statement is the re-graded sphere concentration:

\begin{conjecture}[$Y_1$-A, sphere concentration]
\label{conj:y1a}
For $n\ge3$, $\widetilde H^{k}(\Cn;X)=0$ for $1\le k\le n-2$, and
$\widetilde H^{n-1}(\Cn;X)\cong\sgn_{\Sym_n}$ (one-dimensional). In
particular every non-$\sgn$ isotype of $\widetilde H^{n-1}(\Cn;X)$
vanishes.
\end{conjecture}

This is the re-graded form of the pre-Fork-A target ``UC10.1'': it is
the same spherical-cohomology statement, now stated for the genuine
$\Sym_n$-grading rather than the unconstructed Walsh grading.
$Y_1$-A is \emph{open}. The honest reduction available is the cofiber
long-exact-sequence method, which ports to Fork~A as an
$\Sym_{n-1}$-equivariant reduction --- not, as an earlier draft
claimed, an $\Sym_n$-equivariant one (Theorem~\ref{thm:y1-reduction},
Remark~\ref{rem:y1-branching}).

\begin{construction}[The matched-cylinder prism operator]
\label{constr:bridge}
The doubling functor $\mathrm{db}_x$ embeds $\Cnm$ into $\Cn$ as the
\emph{matched subcategory} $\Dx$, with $X(\mathrm{db}_x\F)\cong
X(\F)\times I$ a prism --- the matched cylinder. The geometric prism
operator $h$, $(h\omega)(D)=\omega(D\times\{0\!\to\!1\})$, is a
degree-$(-1)$ cochain operator. It is a \emph{geometric} construction
on cells; it needs no group action and survives the dropping of
$(\Z/2)^n$ intact. \emph{This much ports cleanly.} What does
\emph{not} port as a discharged input is the twisted-bridge
null-homotopy $\iota_x^{*}\simeq0$ that the cofiber-LES reduction
substitutes; that is the subject of Remark~\ref{rem:bridge-not-ported}
and the sub-residual R1c (Open residual~\ref{res:y1}).
\end{construction}

\begin{remark}[The twisted-bridge null-homotopy does not port as a
discharged input --- sub-residual R1c]
\label{rem:bridge-not-ported}
The pre-Fork-A cofiber-LES reduction collapses its long exact sequence
by substituting the \emph{twisted-bridge null-homotopy}
$\iota_x^{*}\simeq0$. In the pre-Fork-A development
(\texttt{mg-56b8}~\S3.2) that null-homotopy is supplied by the
twisted-bridge chain-homotopy identity
$\iota_x^{*}=\pm\frac12(dh-hd)$, in which the operator is twisted by
the Walsh character $\walsh{\{x\}}$. An earlier draft of this
re-foundation (the text of Construction~\ref{constr:bridge} before
this revision) argued that the identity ``ports without a
$(\Z/2)^n$-action'' because the $\walsh{\{x\}}$-twist is, in Fork~A, a
cochain operation rather than a group action. \textbf{That argument is
a non sequitur, and the claim is withdrawn.} Whether the twist is a
group action or a cochain operation is not what the identity needs:
it needs the cubical coboundary $d$ to \emph{commute} with
multiplication by $\walsh{\{x\}}$. The pre-Fork-A source is explicit
that this very step --- ``$d$ commutes with $\walsh{\{x\}}\cdot$'' ---
is \emph{valid only on $\walsh{S}$-supported cells}
(\texttt{mg-56b8}~\S8.1); so the identity holds, and
$\iota_x^{*}\simeq0$ is established, \emph{only} as a statement about
the $\walsh{S}$-isotype sub-complexes $V_S^{*}$ --- which in the
pre-Fork-A development were genuine sub-complexes, namely the isotypes
of the $(\Z/2)^n$-action on the hocolim.

Fork~A removes exactly that structure. It drops the $(\Z/2)^n$-action
(Status~\ref{status:c-block}), so there is no $(\Z/2)^n$-action on the
Fork~A hocolim and hence \emph{no $V_S^{*}$ sub-complex} for the
null-homotopy to live on; and this document's own
Remark~\ref{rem:walsh-role2} records the underlying fact that $d$ does
\emph{not} preserve harmonic level on the defect complex $X(\F)$ ---
equivalently, $d$ does not commute with $\walsh{\{x\}}$-multiplication
(multiplying by $\walsh{\{x\}}$ shifts Walsh level). Hence in Fork~A
the $\walsh{S}$-graded pieces are \emph{not} sub-complexes, the
twisted-bridge identity does \emph{not} hold as a whole-complex
identity, and the \emph{only} honest pre-Fork-A form of
$\iota_x^{*}\simeq0$ --- ``a map of $\walsh{S}$-isotype sub-complexes
$V_S^{*}$'' --- has \emph{no referent in Fork~A}. The geometric prism
operator $h$ ports; the \emph{twisted} null-homotopy
$\iota_x^{*}\simeq0$ does not port as a discharged input.

This is not a refutation. $\iota_x^{*}\simeq0$ may well be \emph{true}
in the Fork~A setting; it is simply not \emph{discharged} by the
ported twisted bridge. Establishing it in Fork~A --- whether by a
Fork-A-native re-derivation of the null-homotopy, or by carrying the
$\walsh{S}$-grading explicitly and reconciling it with the
$\Sym_{n-1}$-grading of the cofiber LES (a representation-theoretic
reconciliation beside R1b) --- is genuine open work, recorded as the
sub-residual \textbf{R1c} (Open residual~\ref{res:y1}). It is the
third over-claim of the identical ``a $\walsh{S}$-graded fact treated
as portable'' shape that the validate--revise loop has caught: the
\texttt{mg-c15b} self-audit caught it for the $Y_2$
``standard-isotypic'' theorem (Status~\ref{status:y2a}), the
\texttt{mg-d300} validation caught it for the $Y_1$ cofiber-LES
equivariance (R1b), and the \texttt{mg-5221} re-validation caught it
here.
\end{remark}

\begin{theorem}[Cofiber-LES reduction of $Y_1$-A]
\label{thm:y1-reduction}
Fix $n$ and $x\in[n]$. Let $\Cnm\hookrightarrow\Cn$ induce a
cofibration $\iota_x\colon\Xbb_{n-1,x}\hookrightarrow\Xbb_n$ of Fork~A
complexes. Then:
\begin{enumerate}[label=\textup{(\arabic*)},leftmargin=2.4em]
\item The pair $(\Xbb_n,\Xbb_{n-1,x})$ has a cohomology long exact
  sequence equivariant for the point stabiliser
  $\mathrm{Stab}(x)\cong\Sym_{n-1}$ --- and \emph{only} for that
  subgroup, not for all of $\Sym_n$ (Remark~\ref{rem:y1-branching}).
\item Projecting that LES to any $\Sym_{n-1}$-isotype $[\mu]$ (Maschke
  for $\Sym_{n-1}$) and substituting the null-homotopy
  $\iota_x^{*}\simeq0$ --- \emph{when} it is available: in Fork~A it is
  \emph{not} a discharged input but the open sub-residual R1c
  (Construction~\ref{constr:bridge},
  Remark~\ref{rem:bridge-not-ported}) --- collapses it, in each
  degree, to a relation expressing the $[\mu]$-isotype of the
  restriction $\mathrm{Res}^{\Sym_n}_{\Sym_{n-1}}\widetilde
  H^{k}(\Cn;X)$ through the $[\mu]$-isotype cohomology of the cofiber
  $\cofib\iota_x=\Xbb_n/\Xbb_{n-1,x}$ and the smaller-ground-set term
  $\widetilde H^{k-1}(\Xbb_{n-1,x})_{[\mu]}$.
\end{enumerate}
Hence $Y_1$-A reduces to \emph{three} steps: \textup{(a)} control of
the $\Sym_{n-1}$-isotype cohomology of $\cofib\iota_x$, with the
smaller-ground-set terms supplied by the simultaneous induction on
$n$ --- based at the finite check $n=2$ (the families on a two-element
ground set are finite in number and directly computable; this is the
same base verification the pre-Fork-A development relied on for the
$n=2$ sphere class); \textup{(b)} an induction--restriction
(branching) step recovering the $\Sym_n$-isotype structure of
$\widetilde H^{n-1}(\Cn;X)$ --- in particular its
$\sgn_{\Sym_n}$-content, which the $\sgn_{\Sym_n}$ target of $Y_1$-A
(Conjecture~\ref{conj:y1a}) requires --- from the
$\Sym_{n-1}$-equivariant data of step~(a); and \textup{(c)}
establishing the twisted-bridge null-homotopy $\iota_x^{*}\simeq0$ in
the Fork~A setting, which the collapse in~(2) substitutes but which
Construction~\ref{constr:bridge} does \emph{not} discharge
(Remark~\ref{rem:bridge-not-ported}). All three steps are genuinely
open (Open residual~\ref{res:y1}: R1a, R1b, and R1c respectively).
\end{theorem}

\begin{proof}[Proof sketch]
The cofiber long exact sequence of the pair $(\Xbb_n,\Xbb_{n-1,x})$ is
standard. Every map in it is equivariant for the point stabiliser
$\mathrm{Stab}(x)\cong\Sym_{n-1}$ --- the inclusion
$\Cnm\hookrightarrow\Cn$ and the cofiber are natural for relabellings
\emph{that fix $x$} (Remark~\ref{rem:y1-branching}) --- so by Maschke
for $\Sym_{n-1}$ the LES splits into $\Sym_{n-1}$-isotypes. It does
\emph{not} split into $\Sym_n$-isotypes: $\Sym_n$ does not act on the
pair (Remark~\ref{rem:y1-branching}), so Theorem~\ref{thm:maschke},
which is Maschke for $\Sym_n$, does not apply to this LES.
\emph{Granting} the twisted-bridge null-homotopy $\iota_x^{*}\simeq0$
--- which in Fork~A is not a discharged input but the open
sub-residual R1c (Construction~\ref{constr:bridge},
Remark~\ref{rem:bridge-not-ported}) --- substituting the zero map
makes the restriction $j^{*}$ surjective with kernel the image of the
connecting map $\delta$, and $\delta$ injective. This is the
$\Sym_{n-1}$-graded port of the pre-Fork-A cofiber-LES reduction
(\texttt{mg-56b8}~\S3); it carries over an analogous sub-case
structure, including the delicate sub-case where the
smaller-ground-set term meets the sphere class of the $(n-1)$-cube ---
all of it internal to the open cofiber residual R1a. The pre-Fork-A
reduction split into \emph{Walsh} isotypes, for which the
$(\Z/2)^n$-action does descend to the smaller cube; the $\Sym_n$
re-grading has no such descent, which is the source of the branching
step~(b) (Remark~\ref{rem:y1-branching}). Phase~1 does not claim
Fork~A removes the cofiber delicacy, nor that the reduction ports
unchanged: it
claims the reduction ports \emph{to an $\Sym_{n-1}$-equivariant
statement}, and that closing $Y_1$-A then needs all three of the
cofiber (R1a), the branching step (R1b), and a Fork-A-native
discharge of the bridge null-homotopy (R1c).
\end{proof}

\begin{remark}[Why the LES is only $\Sym_{n-1}$-equivariant, and the
branching step this costs]
\label{rem:y1-branching}
The reduction of Theorem~\ref{thm:y1-reduction} acquires, under
Fork~A, a representation-theoretic step the pre-Fork-A reduction did
not have. The smaller object $\Xbb_{n-1,x}$ is the Fork~A complex of
the subcategory $\Cnm\subseteq\Cn$ of pointed families with ground set
inside $[n]\setminus\{x\}$. A relabelling $\pi\in\Sym_n$ carries a
family with ground set $S\subseteq[n]\setminus\{x\}$ to one with
ground set $\pi S\subseteq[n]\setminus\{\pi(x)\}$, so $\pi$ preserves
$\Cnm$ \emph{iff} $\pi(x)=x$: the subcategory $\Cnm$, the inclusion
$\iota_x$, the pair $(\Xbb_n,\Xbb_{n-1,x})$, and its whole cofiber LES
are stable exactly under the point stabiliser
$\mathrm{Stab}(x)\cong\Sym_{n-1}$, not under $\Sym_n$. The
\emph{family} $\{\iota_x\}_{x\in[n]}$ is $\Sym_n$-equivariant (a
permutation carries $\iota_x$ to $\iota_{\pi(x)}$); the
\emph{individual} $\iota_x$, which Theorem~\ref{thm:y1-reduction}
fixes, is only $\mathrm{Stab}(x)$-equivariant.

This is genuinely a cost of dropping $(\Z/2)^n$. The pre-Fork-A
reduction split the cofiber LES into \emph{Walsh} isotypes, and the
$(\Z/2)^n$-action \emph{does} descend to the smaller cube: $(\Z/2)^n$
is abelian and factors as the direct product
$(\Z/2)^{[n]\setminus\{x\}}\times(\Z/2)_x$, so it has a
\emph{quotient} $(\Z/2)^n\twoheadrightarrow(\Z/2)^{[n]\setminus\{x\}}$
acting on $\Xbb_{n-1,x}$ (the $x$-toggle acts trivially --- there is
no $x$-coordinate). $\Sym_n$ has no analogous quotient:
$\mathrm{Stab}(x)\cong\Sym_{n-1}$ is a \emph{subgroup} of $\Sym_n$,
not a quotient, and there is no surjection
$\Sym_n\twoheadrightarrow\Sym_{n-1}$. The abelian product structure
that made the pre-Fork-A cofiber-LES reduction split
isotype-by-isotype is exactly what Fork~A discards when it drops the
$(\Z/2)^n$-action --- so dropping $(\Z/2)^n$ damages the $Y_1$
reduction just as it re-opens the $Y_2$ landing
(Status~\ref{status:c-block}, \S\ref{ssec:y2a}); Phase~1 records both.

\smallskip\noindent\emph{The branching step.} The $\sgn_{\Sym_n}$
target of $Y_1$-A (Conjecture~\ref{conj:y1a}) is an
$\Sym_n$-isotype statement, but Theorem~\ref{thm:y1-reduction}
delivers only $\Sym_{n-1}$-isotype information. Recovering the
$\Sym_n$-isotype content from the $\Sym_{n-1}$-restriction is a
genuine induction--restriction step. It is lossy:
$\sgn_{\Sym_{n-1}}=\mathrm{Res}^{\Sym_n}_{\Sym_{n-1}}\sgn_{\Sym_n}$,
but by the branching rule $\sgn_{\Sym_{n-1}}$ is \emph{also} a
constituent of the restriction of other $\Sym_n$-irreducibles, so an
$\Sym_{n-1}$-isotype computation does not by itself determine the
$\Sym_n$-isotype decomposition --- in particular it does not, by
itself, pin the $\sgn_{\Sym_n}$-content of $\widetilde
H^{n-1}(\Cn;X)$. The branching step is not carried out in Phase~1; it
is a genuine new sub-residual of R1, recorded as R1b (Open
residual~\ref{res:y1}), the $Y_1$-side mirror of the way dropping
$(\Z/2)^n$ re-opens the $Y_2$ landing as R2.
\end{remark}

\begin{openresidual}[The $Y_1$ residual --- R1, in three parts]
\label{res:y1}
$Y_1$-A reduces, by Theorem~\ref{thm:y1-reduction}, to \emph{three}
genuinely open sub-residuals.
\begin{itemize}[leftmargin=1.7em]
\item[\textbf{R1a}] \emph{(the cofiber).} The $\Sym_{n-1}$-isotype
  cohomology of the homotopy cofiber of the matched-cylinder inclusion
  $X(\Dx)\hookrightarrow\Xbb_n$. Stripping the contractible cylinder,
  this is the cohomology of the bar strings of $\Cn$ that \emph{exit}
  the matched subcategory $\Dx$ --- the comma-category twist
  $N(\rho_x/-)$ of the trace functor $\rho_x$, the homotopy type of
  the stuck-family fibres. \emph{As a target object} this is
  unchanged from the pre-Fork-A state: the cofiber-LES re-derivation
  (\texttt{mg-56b8}) had already isolated exactly this object, and
  Fork~A houses it correctly on a well-defined cohomology object. It
  is the genuine open Frankl content and is multi-week-plus work for a
  later phase. The simultaneous induction on $n$ that supplies its
  smaller-ground-set terms is based at the finite check $n=2$.
\item[\textbf{R1b}] \emph{(the branching step --- new under Fork~A).}
  The induction--restriction step recovering the $\Sym_n$-isotype
  structure of $\widetilde H^{n-1}(\Cn;X)$ --- in particular whether
  it is the single $\sgn_{\Sym_n}$ isotype --- from the
  $\Sym_{n-1}$-equivariant cofiber data of R1a
  (Remark~\ref{rem:y1-branching}). This sub-residual \emph{did not
  exist} in the pre-Fork-A development, where the cofiber LES split
  into Walsh isotypes via the abelian-quotient structure of
  $(\Z/2)^n$; it is created by the $(\Z/2)^n\to\Sym_n$ re-grading and
  is the $Y_1$-side mirror of the way dropping $(\Z/2)^n$ re-opens the
  $Y_2$ landing (R2, Open residual~\ref{res:y2}).
\item[\textbf{R1c}] \emph{(the bridge null-homotopy --- not
  discharged under Fork~A).} Establishing the twisted-bridge
  null-homotopy $\iota_x^{*}\simeq0$ as a genuine Fork~A input. The
  collapse of the cofiber LES in Theorem~\ref{thm:y1-reduction}(2)
  substitutes this null-homotopy; the pre-Fork-A development supplied
  it through the twisted-bridge chain-homotopy identity, but
  \emph{only} as a map of the $\walsh{S}$-isotype sub-complexes
  $V_S^{*}$ --- a structure Fork~A does not possess, having dropped
  the $(\Z/2)^n$-action, and on which $d$ does not even preserve the
  harmonic grading (Remark~\ref{rem:walsh-role2}).
  Construction~\ref{constr:bridge} ports the geometric prism operator
  $h$ but \emph{not} the twisted null-homotopy
  (Remark~\ref{rem:bridge-not-ported}). Like R1b, this sub-residual
  \emph{did not exist} in the pre-Fork-A development, where
  $\iota_x^{*}\simeq0$ lived on the genuine $V_S^{*}$ sub-complexes;
  it is created by the $(\Z/2)^n\to\Sym_n$ re-grading. It may be
  closed by a Fork-A-native re-derivation of the null-homotopy, or by
  an explicit $\walsh{S}$-to-$\Sym_{n-1}$ grading reconciliation (a
  second representation-theoretic reconciliation, beside R1b).
\end{itemize}
\textbf{The honest statement of R1 is therefore: its target object
(R1a) is unchanged from the pre-Fork-A state, but its reduction
acquires, under Fork~A, \emph{two} new sub-residuals --- the
$\Sym_{n-1}\to\Sym_n$ branching step R1b and the bridge
null-homotopy R1c --- both created by the $(\Z/2)^n\to\Sym_n$
re-grading.} Fork~A does not close any of the three parts and does
not claim to.
\end{openresidual}

\begin{status}
\label{status:y1a}
$Y_1$-A is \textsf{RED --- open}. This is the expected outcome:
Phase~1 was never going to close the $Y_1$ residual. What Phase~1
establishes is that the cofiber-LES reduction ports to the
well-defined Fork~A object \emph{as an $\Sym_{n-1}$-equivariant
statement} (Theorem~\ref{thm:y1-reduction}), and that the residual is
\emph{three}-part (Open residual~\ref{res:y1}): the cofiber R1a ---
the homotopy type of the stuck-family fibres, the same hard object as
before, now correctly housed --- the new $\Sym_{n-1}\to\Sym_n$
branching sub-residual R1b, and the bridge null-homotopy R1c
($\iota_x^{*}\simeq0$, which the ported twisted bridge does
\emph{not} discharge in Fork~A --- Remark~\ref{rem:bridge-not-ported}).
Both R1b and R1c are created by the $(\Z/2)^n\to\Sym_n$ re-grading.
The pre-Fork-A claim that the reduction ``ports unchanged'' does
\emph{not} hold for Fork~A; the honest statement is that it ports to
an $\Sym_{n-1}$-equivariant reduction \emph{plus} a branching step
\emph{plus} a Fork-A-native discharge of the bridge null-homotopy.
\end{status}

\subsection{The closing contradiction}
\label{ssec:closing}

\begin{theorem}[The Fork A two-part contradiction, conditional]
\label{thm:contradiction}
Assume, for $n\ge3$, the following five inputs, none of which Phase~1
discharges:
\begin{enumerate}[label=\textup{(\roman*)},leftmargin=2.4em]
\item \textup{(R3a, edge-map existence)} the
  \v{C}ech-to-cohomology edge map of Definition~\ref{def:ob}
  \emph{exists} --- the promotion of the mismatch cocycle into the
  cohomology of the category is a defined map, so that $\ob(\Fstar)$
  is a genuine element of $H^{*}(\Cn;X)$;
\item \textup{(R3b, degree behaviour)} that edge map carries the
  mismatch class into total degree $n-1$, so that
  $\ob(\Fstar)\in\widetilde H^{n-1}(\Cn;X)$ specifically;
\item \textup{(R3c, edge-map injectivity)} the edge map is
  \emph{injective} on the class of the mismatch cocycle;
\item \textup{(R2, landing)} $\ob(\Fstar)$ has no
  $\sgn_{\Sym_n}$-component in $\widetilde H^{n-1}(\Cn;X)$ (Open
  residual~\ref{res:y2}); and
\item \textup{(R1, vanishing)} $Y_1$-A holds
  (Conjecture~\ref{conj:y1a}) --- which, by
  Theorem~\ref{thm:y1-reduction}, itself rests on the three-part
  $Y_1$ residual: R1a (the cofiber), R1b (the $\Sym_{n-1}\to\Sym_n$
  branching step), and R1c (the twisted-bridge null-homotopy
  $\iota_x^{*}\simeq0$, not discharged by the ported bridge ---
  Remark~\ref{rem:bridge-not-ported}), Open residual~\ref{res:y1}.
\end{enumerate}
Then there is no minimal counterexample to Frankl on $[n]$; hence
Frankl holds on $[n]$.
\end{theorem}

\begin{proof}
Suppose $\Fstar$ is a minimal counterexample. By (i) and (ii) the
obstruction class $\ob(\Fstar)$ is a well-defined element of
$\widetilde H^{n-1}(\Cn;X)$ --- (i) makes the promotion a defined map,
(ii) places its value in degree $n-1$ --- so that Fact One and the
remainder of the argument have a subject. By Fact One
(Proposition~\ref{prop:fact-one}, using the injectivity (iii)),
$\ob(\Fstar)\neq0$ in $\widetilde H^{n-1}(\Cn;X)$. By (v), $Y_1$-A
makes $\widetilde H^{n-1}(\Cn;X)\cong\sgn_{\Sym_n}$ --- the cohomology
in degree $n-1$ is exactly the one-dimensional sphere class, all other
isotypes vanish. By (iv), $\ob(\Fstar)$ has no $\sgn$-component; since
$\widetilde H^{n-1}(\Cn;X)$ is \emph{entirely} $\sgn$-isotypic, the
only element with no $\sgn$-component is $0$, so $\ob(\Fstar)=0$. This
contradicts $\ob(\Fstar)\neq0$. Hence no minimal counterexample
exists, and Frankl holds on $[n]$.
\end{proof}

\begin{remark}[The honest shape of the re-founded contradiction]
\label{rem:honest-shape}
Theorem~\ref{thm:contradiction} is the Fork~A re-founding of the
intended proof. It is a \emph{conditional} theorem, and the honest
accounting of its hypotheses is the heart of the Phase-1 deliverable.
\begin{itemize}[leftmargin=1.5em]
\item It rests on the named open inputs R1, R2, R3 --- with R1 itself
  three-part (R1a the cofiber, R1b the $\Sym_{n-1}\to\Sym_n$ branching
  step, R1c the twisted-bridge null-homotopy) and R3 itself three-part
  (R3a existence, R3b degree, R3c injectivity of the edge map).
  Phase~1 makes \emph{no} claim that this is a count of ``exactly
  three'': it is the honest enumeration of every input
  Theorem~\ref{thm:contradiction} rests on. The comparison with the
  pre-Fork-A closing argument is genuine but modest: Fork~A removes
  the dependence on the unconstructed $(\Z/2)^n$-action
  (Status~\ref{status:c-block}) and houses the construction on a
  well-defined object --- a gain in soundness and clarity --- but in
  doing so it trades the $(\Z/2)^n$ input for the two $Y_1$-side
  sub-residuals R1b and R1c, and surfaces the edge-map
  existence/degree debt R3a/R3b. It is therefore \emph{not} a clean
  reduction in the number of open inputs, and Phase~1 does not claim
  one.
\item R1a (the $Y_1$ cofiber residual, as a target object) and R3c
  (edge-map injectivity) are \emph{unchanged} from the pre-Fork-A
  state. The inputs Fork~A genuinely changes are four: R2 (the
  landing), which it re-opens --- the pre-Fork-A $Y_2$ was
  ``mechanism-sound'' resting on $(\Z/2)^n$-equivariance, and dropping
  that action means the separation of $\ob$ from the sphere class must
  be re-established by other means (Open residual~\ref{res:y2}); R1b
  (the $\Sym_{n-1}\to\Sym_n$ branching step), which the
  $(\Z/2)^n\to\Sym_n$ re-grading newly creates
  (Remark~\ref{rem:y1-branching}); R1c (the twisted-bridge
  null-homotopy $\iota_x^{*}\simeq0$), which the same re-grading
  surfaces --- the ported twisted bridge established it only on the
  $\walsh{S}$-isotype sub-complexes $V_S^{*}$, which Fork~A, having
  dropped $(\Z/2)^n$, does not possess
  (Remark~\ref{rem:bridge-not-ported}); and R3a/R3b (edge-map
  existence and degree), implicit and unexamined in the pre-Fork-A
  development's undefined object, which Fork~A must now name because
  its object is genuinely defined.
\item The proof of Theorem~\ref{thm:contradiction} does not invoke
  spherical concentration as an independent ingredient beyond $Y_1$-A;
  spherical concentration \emph{is} $Y_1$-A. It re-enters as the
  inductive partner of $Y_1$-A in the cofiber-LES method
  (Theorem~\ref{thm:y1-reduction}), exactly as in the pre-Fork-A
  simultaneous induction.
\end{itemize}
None of this proves Frankl: the inputs (i)--(v) of
Theorem~\ref{thm:contradiction} are genuinely open. What
Theorem~\ref{thm:contradiction} shows is that the re-founded
contradiction is coherent on a sound object, and that its open part is
named precisely and \emph{in full} --- the residuals R1 (three-part),
R2, and R3 (three-part).
\end{remark}

% =====================================================================
\section{The two owed debts}
\label{sec:debts}
% =====================================================================

\subsection{Owed debt 1 --- the re-grading: ADDRESSED}
\label{ssec:debt1}

\begin{status}
\label{status:debt1}
Owed debt~\ref{debt:regrading} (the re-grading) is \textsf{ADDRESSED}:
its structural content is discharged, and its one remaining open part
is precisely localized.
\end{status}

\noindent\emph{Discharged (structural re-grading).}
\begin{enumerate}[label=\textup{(\arabic*)},leftmargin=2.2em]
\item The $(\Z/2)^n$-action is dropped, which \emph{dissolves}
  Phase-1.5 refutation (b) (Status~\ref{status:c-block}).
\item The grading is re-established by Maschke for the genuine
  $\Sym_n$-action (Theorem~\ref{thm:maschke}) --- the genuine analogue
  of the old Walsh decomposition, for the group that genuinely acts.
\item The sphere class is identified as the $\sgn_{\Sym_n}$ isotype
  (Proposition~\ref{prop:sphere-sgn}); the old target's $\walsh{[n]}$
  factor was redundant $(\Z/2)^n$-decoration
  (Remark~\ref{rem:sgn-always-sn}).
\item The vanishing statement $Y_1$ re-derives to the re-graded
  sphere concentration $Y_1$-A (Conjecture~\ref{conj:y1a}), with the
  cofiber-LES reduction ported to an $\Sym_{n-1}$-equivariant
  statement (Theorem~\ref{thm:y1-reduction}).
\item The landing statement $Y_2$ re-derives at the level of its
  structural mechanism (Proposition~\ref{prop:y2a-mech}).
\end{enumerate}

\noindent\emph{Localized (the open parts).} The re-grading correctly
\emph{surfaces} three inputs that the pre-Fork-A development obtained
for free from the $(\Z/2)^n$-action now dropped.
\begin{enumerate}[label=\textup{(\roman*)},leftmargin=2.2em]
\item \emph{The $Y_2$ separation.} The separation of $\ob$ from the
  sphere class, free in the Walsh grading, is not free in the $\Sym_n$
  grading. Phase~1 localizes this precisely as residual R2 (Open
  residual~\ref{res:y2}).
\item \emph{The $Y_1$ branching step.} The pre-Fork-A cofiber-LES
  reduction split into Walsh isotypes, descending to the smaller cube
  via the abelian-quotient structure of $(\Z/2)^n$; $\Sym_n$ has no
  such quotient (Remark~\ref{rem:y1-branching}). The reduction
  therefore ports only to an $\Sym_{n-1}$-equivariant statement, and
  recovering the $\Sym_n$-isotype structure needs an
  $\Sym_{n-1}\to\Sym_n$ branching step. Phase~1 localizes this
  precisely as the sub-residual R1b (Open residual~\ref{res:y1}).
\item \emph{The $Y_1$ bridge null-homotopy.} The pre-Fork-A
  cofiber-LES reduction collapses its LES by substituting the
  twisted-bridge null-homotopy $\iota_x^{*}\simeq0$, which the
  pre-Fork-A twisted bridge established only on the
  $\walsh{S}$-isotype sub-complexes $V_S^{*}$. Fork~A, having dropped
  the $(\Z/2)^n$-action, has no $V_S^{*}$ sub-complex, so the ported
  bridge does not discharge $\iota_x^{*}\simeq0$
  (Remark~\ref{rem:bridge-not-ported}). Phase~1 localizes this
  precisely as the sub-residual R1c (Open residual~\ref{res:y1}).
\end{enumerate}
None of the three is a failure of the re-grading; all three are the
re-grading doing its job --- exposing that inputs the pre-Fork-A
ledger treated as settled (``$Y_2$ mechanism-sound''; ``$Y_1$
reduction unchanged'') in fact depended on the now-refuted
$(\Z/2)^n$-action and must be re-established. The discipline is
symmetric: dropping $(\Z/2)^n$ re-opens work on \emph{both} sides of
the contradiction --- the $Y_2$ landing (R2) on one side, the $Y_1$
branching step (R1b) and bridge null-homotopy (R1c) on the other ---
and Phase~1 records them all.

\begin{remark}[Why ``ADDRESSED'' and not ``PAID'']
\label{rem:debt1-addressed}
An earlier draft of this document recorded Owed debt~\ref{debt:regrading}
as fully \textsf{PAID}, on the strength of a claimed theorem that
$\ob$ is standard-isotypic. That theorem does not hold (Status~%
\ref{status:y2a}), so the honest status is \textsf{ADDRESSED}: the
re-grading as a structural task is complete and sound, and its
residue is precisely named --- the landing residual R2 and the two
$Y_1$-side sub-residuals R1b (branching) and R1c (the bridge
null-homotopy) that the re-grading surfaces. The brief
permits ``address \emph{or} precisely localize''; Phase~1 does both
--- addresses the structural re-grading, localizes the residue.
\end{remark}

\subsection{Owed debt 2 --- faithfulness ($V_3$ / wrapper-risk): LOCALIZED}
\label{ssec:debt2}

\begin{status}
\label{status:debt2}
Owed debt~\ref{debt:faithful} (faithfulness) is \textsf{LOCALIZED, not
discharged}. Phase~1 reduces it to a single named edge map and three
load-bearing properties of that map --- existence, degree-$(n-1)$
landing, and injectivity (R3a, R3b, R3c) --- with one further facet
(R2) surfaced by the re-grading.
\end{status}

The vision's clause $V_3$ --- ``the minimal counterexample differs
from the generic families by a factor relating to this cohomology''
--- asks that the cohomology \emph{genuinely detect} the
counterexample. On the Fork~A object the ``factor'' is the obstruction
class $\ob(\Fstar)$: a generic family (one with a rare coordinate) has
$\ob=0$, because its local rare-coordinate data glue through its own
rare coordinate; a minimal counterexample has $\ob\neq0$ ---
\emph{provided} Fact One holds. So $\ob$ distinguishes $\Fstar$ from
the generic families, as $V_3$ asks, conditionally on Fact One.

\begin{debt}[Faithfulness, localized --- R3, a three-part debt]
\label{debt:faithful-localized}
Owed debt~\ref{debt:faithful} is the conjunction of \emph{three}
statements about the \v{C}ech-to-cohomology edge map of
Definition~\ref{def:ob} --- all three load-bearing, none discharged by
Phase~1.
\begin{itemize}[leftmargin=1.7em]
\item[\textbf{R3a}] \emph{(existence).} The promotion of the mismatch
  cocycle through the \v{C}ech-to-cohomology spectral sequence into
  $H^{*}(\Cn;X)$ is a \emph{defined map}. This is not automatic: the
  \v{C}ech double complex $\check C^{*,*}(\mathfrak U;\fobs)$ lives on
  the punctured trace poset $\Cn(\Fstar)\setminus\{\Fstar\}$, while
  $H^{*}(\Cn;X)$ is the cohomology of the whole category $\Cn$, and
  the comparison map between the two is not constructed in Phase~1
  (Remark~\ref{rem:promotion}).
\item[\textbf{R3b}] \emph{(degree-$(n-1)$ landing).} The edge map
  carries the mismatch class --- which sits in \v{C}ech degree $1$ ---
  into total degree $n-1$, so that $\ob(\Fstar)$ is an element of
  $\widetilde H^{n-1}(\Cn;X)$ \emph{specifically}, the degree where
  the sphere-concentration residual R1 applies to it. Phase~1 inherits
  the degree count $n-1$ from the pre-Fork-A development
  (Remark~\ref{rem:promotion}), where it was performed in the
  now-defunct object $X^{\cap}_n$; it is not re-verified in the
  genuinely different Fork~A object.
\item[\textbf{R3c}] \emph{(injectivity).} The edge map $\check
  H^{1}(\mathfrak U;\fobs)\to\widetilde H^{n-1}(\Cn;X)$ is
  \emph{injective} on the class of the mismatch cocycle: if
  $\ob(\Fstar)=0$ in $\widetilde H^{n-1}(\Cn;X)$, then $\{m_{xy}\}$ is
  the \v{C}ech coboundary of a \emph{global} section of $\fobs$.
\end{itemize}
R3c is the original faithfulness statement: the promotion of the
obstruction from the \v{C}ech $H^1$ of the cover into the cohomology
of the category loses no information --- a vanishing promoted class
reflects a genuinely trivial mismatch. R3a and R3b are its logical
prerequisites: without them $\ob(\Fstar)$ is not a well-defined
element of $\widetilde H^{n-1}(\Cn;X)$ at all, and R3c --- and Fact
One, and Theorem~\ref{thm:contradiction} --- have no subject. Phase~1
names all three and discharges none. \emph{The earlier Phase-1 draft
named only R3c}; that the residual set is a three-part debt, not a
single injectivity statement, is one of the two corrections of the
\texttt{mg-365e} revision (Status~\ref{status:mg365e-revision}).
\end{debt}

Three honest remarks sharpen what this debt is and how it relates to
the rest of the contradiction.

\begin{remark}[The wrapper-risk, precisely]
\label{rem:wrapper}
The standing concern (``wrapper-risk'') is that the vanishing side of
the contradiction --- $Y_1$-A --- uses only the shape of the \v{C}ech
double complex and never the counterexample hypothesis, so that
\emph{if} $Y_1$-A holds then $\ob(\F)=0$ for \emph{every} family, and
all Frankl-specific content collapses into Fact One. Fork~A does not
remove this risk, but it localizes it: under Fork~A the
Frankl-specific content is the edge-map injectivity R3c. If that
injectivity holds, the equivalence ``$\ob=0\iff$ a global
rare-coordinate section exists'' is genuine and the cohomology
genuinely detects the counterexample ($V_3$ satisfied). If it fails,
$\ob$ is family-blind and the contradiction is hollow. The debt is a
small, checkable package of homological-algebra statements about one
edge map --- its existence, its degree, and its injectivity
(R3a--R3c).
\end{remark}

\begin{remark}[R2 is faithfulness-flavoured]
\label{rem:r2-faithful}
The landing residual R2 (Open residual~\ref{res:y2}) --- whether
$\ob$ avoids the $\sgn$ sphere-class isotype --- is closely related to
faithfulness. If $\ob$ had a $\sgn$-component, that component would be
a multiple of the sphere class, which is family-\emph{independent}
(it is a property of the category's cohomology, not of any family). An
$\ob$ with a nonzero $\sgn$-component would thus be partly a
family-blind invariant --- a faithfulness failure of exactly the
wrapper-risk shape. So R2 and R3, though housed under the two
different debts, are facets of one underlying question: whether the
cohomology separates the family-specific obstruction from the
family-blind orientation class. Phase~1 records them as two named
residuals; a later phase may find them inter-reducible.
\end{remark}

\begin{remark}[The hocolim-versus-holim sharpening; the foundational prerequisite]
\label{rem:holim-prereq}
Remark~\ref{rem:hocolim-flag} flagged that the obstruction is a
gluing-type (limit-flavoured) invariant while $H^{*}(\Cn;X)$ is the
cohomology of a homotopy \emph{colimit}. This is a sharpening of R3,
not a separate debt: the edge map of Definition~\ref{def:ob} is the
comparison between the limit-side \v{C}ech $H^1$ (where gluing lives)
and the colimit-side $H^{n-1}(\Cn;X)$ (where $\ob$ is housed), and
edge-map injectivity is exactly the assertion that this comparison is
faithful. Pinning the edge map to the precision R3 needs requires the
cell-level Bousfield--Kan model of $\Cn$: the projection $\pi_T$ at
cell level (supplied --- Theorem~\ref{thm:projection}) and the
cell-level structure of $\fobs$ and its spectral sequence (not pinned
here). Fork~A discharges the harder half of this foundational
prerequisite (the structure map genuinely exists and is pinned) and
leaves the spectral sequence of $\fobs$ for the residual-attack phase.
A computation reported on un-pinned $\fobs$-foundations would repeat
the Phase-1.5 failure mode and must not be attempted before the
pinning.
\end{remark}

% =====================================================================
\section{Verdict, what Phase 1 delivered, recommendations}
\label{sec:verdict}
% =====================================================================

\subsection{Verdict}
\label{ssec:verdict}

\textbf{GREEN --- FORK-A-REFOUNDATION-DELIVERED (Phase~1).} The Frankl
cohomology is re-founded on Fork~A in LaTeX. The covariant homotopy
colimit $H^{*}(\Cn;X)$ is a genuinely defined cohomology theory
(Theorem~\ref{thm:welldefined}), built on the trace projection
$\pi_T$, which exists and is functorial (Theorem~\ref{thm:projection}):
this dodges Phase-1.5 refutation (a) by construction. The
$(\Z/2)^n$-action is dropped, dissolving Phase-1.5 refutation (b)
(Status~\ref{status:c-block}). The grading is re-established by Maschke
for the genuine $\Sym_n$-action (Theorem~\ref{thm:maschke}), with the
sphere class identified as $\sgn_{\Sym_n}$
(Proposition~\ref{prop:sphere-sgn}). The obstruction class and the
two-part contradiction are re-established as a conditional theorem
(Theorem~\ref{thm:contradiction}). Owed debt~\ref{debt:regrading} (the
re-grading) is addressed --- its structural content discharged, its
residue localized as R2 (the $Y_2$ separation) and the two
$Y_1$-side sub-residuals the re-grading newly surfaces, R1b (the
$\Sym_{n-1}\to\Sym_n$ branching step) and R1c (the twisted-bridge
null-homotopy); Owed debt~\ref{debt:faithful} (faithfulness) is
localized as the three-part R3. \textbf{Frankl is not proved, and
Phase~1 was never meant to prove it.} The verdict GREEN is for the
Phase-1 scope: the re-foundation is delivered and internally sound.
The genuine open content is named precisely and in full --- R1
(three-part: R1a the $Y_1$ cofiber, R1b the branching step, R1c the
twisted-bridge null-homotopy), R2 (the $Y_2$ landing / separation),
and R3 (three-part: existence, degree, and injectivity of the
\v{C}ech-to-cohomology edge map).

\subsection{What Phase 1 delivered (positive content)}
\label{ssec:delivered}

\begin{enumerate}[leftmargin=1.9em]
\item A genuinely defined cohomology object $H^{*}(\Cn;X)$ --- the
  covariant hocolim, with $D^2=0$ verified
  (Theorem~\ref{thm:welldefined}) --- replacing the undefined
  pre-Fork-A $X^{\cap}_n$.
\item The trace projection $\pi_T$, pinned as a covariant functor
  (Theorem~\ref{thm:projection}), the structure map the build had
  backwards.
\item The re-grading: $\Sym_n$-isotype decomposition by Maschke
  (Theorem~\ref{thm:maschke}); the sphere class as $\sgn_{\Sym_n}$
  (Proposition~\ref{prop:sphere-sgn}); the $(\Z/2)^n$-action dropped
  with no loss at the sphere class.
\item The obstruction class re-founded on the well-defined object
  (\S\ref{sec:obstruction}), and the structural mechanism of the
  landing re-derived in the $\Sym_n$ grading with no group action
  (Proposition~\ref{prop:y2a-mech}).
\item The $Y_1$ vanishing re-stated as re-graded sphere concentration
  (Conjecture~\ref{conj:y1a}), with the cofiber-LES reduction ported
  to an $\Sym_{n-1}$-equivariant statement
  (Theorem~\ref{thm:y1-reduction}) and the two new sub-residuals it
  surfaces honestly recorded --- the $\Sym_{n-1}\to\Sym_n$ branching
  step R1b (Remark~\ref{rem:y1-branching}) and the twisted-bridge
  null-homotopy R1c (Remark~\ref{rem:bridge-not-ported}).
\item The closing contradiction re-established as a conditional
  theorem (Theorem~\ref{thm:contradiction}), with \emph{every} input
  it rests on enumerated in full --- R1 (three-part: R1a cofiber, R1b
  branching, R1c the bridge null-homotopy), R2, and R3 (three-part:
  R3a existence, R3b degree, R3c injectivity of the edge map) --- and
  no claim of a clean ``exactly three / one fewer than pre-Fork-A''
  count (Remark~\ref{rem:honest-shape}).
\item Owed debt~1 addressed; Owed debt~2 localized
  (\S\ref{sec:debts}); and the over-claims of the
  ``$\walsh{S}$-graded fact treated as portable'' shape caught and
  withdrawn --- the ``standard-isotypic'' landing theorem
  (Status~\ref{status:y2a}, Remark~\ref{rem:debt1-addressed}) and the
  twisted-bridge null-homotopy (Construction~\ref{constr:bridge},
  Remark~\ref{rem:bridge-not-ported}, now the sub-residual R1c).
\end{enumerate}

\subsection{What remains open}
\label{ssec:open}

\begin{itemize}[leftmargin=1.7em]
\item \textbf{R1 --- $Y_1$-A, Open residual~\ref{res:y1}}, in three
  parts. \textbf{R1a:} the $\Sym_{n-1}$-isotype cohomology of the
  cofiber of the matched-cylinder inclusion --- the homotopy type of
  the stuck-family fibres; as a target object unchanged from the
  pre-Fork-A residual, the genuine open Frankl content,
  multi-week-plus. \textbf{R1b:} the $\Sym_{n-1}\to\Sym_n$ branching
  step recovering the $\Sym_n$-isotype structure of $\widetilde
  H^{n-1}(\Cn;X)$ from the $\Sym_{n-1}$-equivariant cofiber data --- a
  new sub-residual created by the $(\Z/2)^n\to\Sym_n$ re-grading
  (Remark~\ref{rem:y1-branching}). \textbf{R1c:} the twisted-bridge
  null-homotopy $\iota_x^{*}\simeq0$, which the cofiber-LES collapse
  substitutes but the ported bridge does not discharge in Fork~A ---
  also new, created by the same re-grading
  (Remark~\ref{rem:bridge-not-ported}).
\item \textbf{R2 --- $Y_2$-A, Open residual~\ref{res:y2}.} Whether
  $\ob(\Fstar)$ avoids the $\sgn$ sphere-class isotype. A
  representation-theoretic question about the obstruction sheaf
  $\fobs$; faithfulness-flavoured (Remark~\ref{rem:r2-faithful}).
\item \textbf{R3 --- faithfulness, Owed
  debt~\ref{debt:faithful-localized}}, a three-part debt on the
  \v{C}ech-to-cohomology edge map: \textbf{R3a} its existence,
  \textbf{R3b} its degree-$(n-1)$ landing, \textbf{R3c} its
  injectivity --- together the whole of the $V_3$/wrapper-risk
  question; subsumes the hocolim-versus-holim question
  (Remark~\ref{rem:holim-prereq}).
\end{itemize}

\subsection{Recommendations (routed to Daniel via \texttt{pm-onethird})}
\label{ssec:recommendations}

\begin{enumerate}[leftmargin=1.9em]
\item \textbf{Proceed to a further independent re-validation of this
  twice-revised document.} Two Phase-2 validations have now been
  incorporated: \texttt{mg-d300} (AMBER) returned two over-claims,
  corrected by the \texttt{mg-365e} revision
  (Status~\ref{status:mg365e-revision}); \texttt{mg-5221} (AMBER)
  returned a third of the same shape, corrected by the
  \texttt{mg-894c} revision (Status~\ref{status:mg894c-revision}).
  The next re-validation should default-skeptically confirm that the
  correction is complete and honest --- in particular that
  Construction~\ref{constr:bridge} no longer claims the twisted
  bridge ``ports'' and that R1c is genuinely recorded
  (Remark~\ref{rem:bridge-not-ported}, Open residual~\ref{res:y1});
  that Theorem~\ref{thm:contradiction} lists \emph{every} input, with
  R1 now three-part and R3 three-part; that the exhaustive Walsh-fact
  sweep recorded in \texttt{docs/Frankl-ForkA-Phase2-Revision2.md} is
  complete --- no fourth ported $\walsh{S}$-graded fact treated as
  portable; and that the spine --- Theorem~\ref{thm:projection}
  (chain map and functoriality), Theorem~\ref{thm:welldefined}
  ($D^2=0$), Theorem~\ref{thm:maschke} (the $\Sym_n$-equivariance of
  $D$), Proposition~\ref{prop:y2a-mech} --- which \texttt{mg-d300}
  and \texttt{mg-5221} both independently confirmed sound, is
  undisturbed by the revisions.
\item \textbf{The $Y_1$ residual} (R1, Open residual~\ref{res:y1}) and
  \textbf{the faithfulness edge map} (R3, Owed
  debt~\ref{debt:faithful-localized}) can be attacked in parallel as
  separate tickets, with the corrected picture: R1 now carries,
  besides the cofiber R1a, \emph{two} added sub-residuals the
  re-grading surfaces --- the $\Sym_{n-1}\to\Sym_n$ branching step
  R1b (Remark~\ref{rem:y1-branching}) and the twisted-bridge
  null-homotopy R1c (Remark~\ref{rem:bridge-not-ported}); and R3 is
  three sub-problems --- existence R3a, degree R3b, injectivity R3c
  --- of which existence and degree are upstream of injectivity. Both
  attacks require first pinning the cell-level spectral sequence of
  $\fobs$ (Remark~\ref{rem:holim-prereq}).
\item \textbf{The $Y_2$ separation residual} (R2, Open
  residual~\ref{res:y2}) should be attacked together with R3: they are
  facets of one faithfulness question (Remark~\ref{rem:r2-faithful}),
  and a representation-theoretic analysis of $\fobs$ feeds both.
\item \textbf{Do not Lean-formalize yet} (Phase~3). Phase~3 is gated on
  a further independent re-validation of this twice-revised document
  returning GREEN. Of the Fork~A content,
  Theorems~\ref{thm:projection},~\ref{thm:welldefined}, and~%
  \ref{thm:maschke} and Proposition~\ref{prop:y2a-mech} are the
  formalization-ready core; the residuals R1 (three-part), R2, and R3
  (three-part) are not.
\item \textbf{The consolidated paper (\texttt{../main.tex}) is not
  edited by Phase~1.} It honestly records the pre-Fork-A state. A
  later pass --- once a further independent re-validation of this
  twice-revised document returns GREEN --- may fold the Fork~A
  re-foundation into the paper, superseding its \S\S3--4. That is out
  of Phase-1 scope.
\end{enumerate}

\subsection{The honest one-paragraph summary}
\label{ssec:final}

Fork~A re-founds the Frankl cohomology as the covariant homotopy
colimit of the per-family cubical-defect diagram over the trace
category, with the canonical trace projection --- which genuinely
exists and is functorial --- as its structure map; this object is
well-defined, which the pre-Fork-A object was not, so the Phase-1.5
against-trace refutation is dodged by construction. Dropping the
$(\Z/2)^n$-action dissolves the Phase-1.5 $(\Z/2)^n$ refutation
outright, because there is then no group action to fail, and the
grading is re-established by Maschke for the genuine $\Sym_n$-action,
with the sphere class the sign isotype --- the sign having always been
an $\Sym_n$-datum, not a $(\Z/2)^n$ one. On the re-founded object the
obstruction class is re-established and the two-part contradiction is
re-assembled as a conditional theorem; honest accounting enumerates,
in full, every input it rests on --- the $Y_1$ cofiber residual R1
(its target object R1a unchanged, but its reduction now carrying two
new sub-residuals the $(\Z/2)^n\to\Sym_n$ re-grading creates: the
$\Sym_{n-1}\to\Sym_n$ branching step R1b, because the pre-Fork-A
reduction split isotype-by-isotype only via the abelian-quotient
structure of $(\Z/2)^n$ that $\Sym_n$ lacks; and the twisted-bridge
null-homotopy R1c, because the pre-Fork-A bridge established
$\iota_x^{*}\simeq0$ only on the $\walsh{S}$-isotype sub-complexes
$V_S^{*}$ that Fork~A, having dropped $(\Z/2)^n$, no longer has), the
$Y_2$ landing R2
(which Fork~A genuinely re-opens, because the pre-Fork-A landing
mechanism rested on the now-dropped $(\Z/2)^n$-equivariance, and the
separation of the obstruction from the sign sphere class is a
representation-theoretic question about the obstruction sheaf that
Phase~1 localizes but does not settle), and a three-part faithfulness
debt R3 --- the existence, the degree-$(n-1)$ behaviour, and the
injectivity of the \v{C}ech-to-cohomology edge map that promotes the
obstruction into the cohomology (the vision's clause $V_3$). Frankl is not proved, and Phase~1
was not meant to prove it: what Phase~1 delivers is a sound
re-foundation that clears both Phase-1.5 walls, carries out the
re-grading, re-assembles the contradiction on a well-defined object,
and isolates the genuine open content as the precisely and completely
named residuals R1 (three-part), R2, and R3 (three-part) --- and, in
the discipline the ledger's over-claim history demands, it withdraws
over-claims rather than carrying them forward: the
clean-landing-theorem over-claim caught in Phase~1, the two caught by
the \texttt{mg-d300} Phase-2 validation and corrected in the
\texttt{mg-365e} revision (the ``$\Sym_n$-equivariant cofiber LES /
$R1$ unchanged'' over-claim and the ``exactly three residuals''
over-claim; Status~\ref{status:mg365e-revision}), and the third of
the identical ``$\walsh{S}$-graded fact treated as portable'' shape
caught by the \texttt{mg-5221} Phase-2 re-validation and corrected in
the \texttt{mg-894c} revision (the twisted-bridge null-homotopy
$\iota_x^{*}\simeq0$, presented as a discharged ported input, now
the sub-residual R1c; Status~\ref{status:mg894c-revision}). After the
\texttt{mg-894c} exhaustive Walsh-fact sweep, the input list of
Theorem~\ref{thm:contradiction} is complete and final.

\subsection{Cross-references}
\label{ssec:xref}

\begin{itemize}[leftmargin=1.7em]
\item \textbf{Brief:} \texttt{mg-c15b} (Frankl-ForkA-Refound-LaTeX,
  Phase~1); \texttt{mg-365e} (Frankl-ForkA-Phase2-Revision);
  \texttt{mg-894c} (Frankl-ForkA-Phase2-Revision2 --- the revision
  carried out in this document).
\item \textbf{The Phase-2 independent validations (the findings the
  two revisions correct):}
  \texttt{docs/Frankl-ForkA-Phase2-validation.md} (\texttt{mg-d300};
  \S3 Finding~A, \S4 Finding~B, \S5.1 the minor point, \S6.2 the
  revision spec) and
  \texttt{docs/Frankl-ForkA-Phase2-Revalidation.md} (\texttt{mg-5221};
  \S3 the surviving twisted-bridge over-claim, \S4 the minor points,
  \S5.2 the re-statement spec); the revisions are logged in
  \texttt{docs/Frankl-ForkA-Phase2-Revision.md} (\texttt{mg-365e})
  and \texttt{docs/Frankl-ForkA-Phase2-Revision2.md}
  (\texttt{mg-894c} --- the latter also carrying the exhaustive
  Walsh-fact sweep).
\item \textbf{Fork~A endorsement and the two debts:}
  \texttt{docs/Frankl-ForkD-Probe-n3.md} (\texttt{mg-565a}, \S\S6--7).
\item \textbf{The drift diagnosis and the fork menu:}
  \texttt{docs/Frankl-vision-check.md} (\texttt{mg-6edc}, \S\S6--7).
\item \textbf{The Phase-1.5 refutations and the pinned projection
  $\pi_T$:} \texttt{docs/Frankl-Y1-phase1.5-bk-structure-maps.md}
  (\texttt{mg-0405}: \S3 refutation (a), \S4 refutation (b), \S3.6
  Proposition~3.6 = Theorem~\ref{thm:projection} here, \S6 Fork~A).
\item \textbf{The cofiber-LES reduction of $Y_1$ and its residual:}
  \texttt{docs/Frankl-Y1-reprove.md} (\texttt{mg-56b8}).
\item \textbf{The consolidated paper (pre-Fork-A state):}
  \texttt{../main.tex} and \texttt{../sections/03--06}.
\item \textbf{Onboarding:} \texttt{docs/PROOF-STRUCTURE-ONBOARDING.md}.
\end{itemize}

\bigskip
\noindent\rule{\textwidth}{0.4pt}

\medskip
\noindent\footnotesize Phase~1 re-foundation by polecat
\texttt{cat-mg-c15b}, revised by polecat \texttt{cat-mg-365e} per the
\texttt{mg-d300} Phase-2 validation and by polecat
\texttt{cat-mg-894c} per the \texttt{mg-5221} Phase-2 re-validation.
Deliverable on \texttt{main}: this document
(\texttt{paper/forkA/refoundation.tex}) and the ledger notes
\texttt{docs/Frankl-ForkA-Refound-LaTeX.md} (Phase~1),
\texttt{docs/Frankl-ForkA-Phase2-Revision.md} (the \texttt{mg-365e}
revision), and \texttt{docs/Frankl-ForkA-Phase2-Revision2.md} (the
\texttt{mg-894c} revision).
Verdict: \textbf{GREEN --- FORK-A-REFOUNDATION-DELIVERED (Phase~1),
as corrected.} The Frankl cohomology is re-founded on Fork~A: the
covariant homotopy colimit $H^{*}(\Cn;X)$ is genuinely defined
(dodging Phase-1.5 refutation (a) by construction), the
$(\Z/2)^n$-action is dropped (dissolving refutation (b)), the
$\sgn_{\Sym_n}$ re-grading is carried out, and the obstruction class
and two-part contradiction are re-established as a conditional theorem
whose open inputs are named precisely and in full --- R1 ($Y_1$
cofiber: R1a unchanged as a target object, R1b a new
$\Sym_{n-1}\to\Sym_n$ branching step, R1c the twisted-bridge
null-homotopy the ported bridge does not discharge), R2 ($Y_2$
landing, re-opened and localized), and the three-part R3 (edge-map
existence, degree, and injectivity). Owed debt~1 (re-grading) is
addressed; Owed debt~2 (faithfulness) is localized. Frankl is not
proved; Phase~1 was not meant to prove it. \emph{Revision note
(\texttt{mg-365e}, \texttt{mg-894c}).} This document has been revised
twice following independent Phase-2 validation. The \texttt{mg-365e}
revision acted on the \texttt{mg-d300} validation (AMBER), which
caught two over-claims: that the $Y_1$ cofiber-LES reduction is
$\Sym_n$-equivariant and ``ports unchanged'' (it is only
$\Sym_{n-1}=\mathrm{Stab}(x)$-equivariant, and acquires the
$\Sym_{n-1}\to\Sym_n$ branching sub-residual R1b), and that the
contradiction rests on ``exactly three'' residuals (R3 is in fact a
three-part edge-map debt). The \texttt{mg-894c} revision acted on the
\texttt{mg-5221} re-validation (AMBER), which caught a third
over-claim of the identical ``$\walsh{S}$-graded fact treated as
portable'' shape: the twisted-bridge null-homotopy
$\iota_x^{*}\simeq0$ was presented as a discharged, ported input
(Construction~\ref{constr:bridge}) when it is a
$\walsh{S}$-isotype-sub-complex fact with no Fork~A referent --- now
the sub-residual R1c (Remark~\ref{rem:bridge-not-ported}); the
\texttt{mg-894c} revision additionally carried out an exhaustive
sweep of every ported Walsh fact (ledger
\texttt{docs/Frankl-ForkA-Phase2-Revision2.md}), certifying there is
no fourth prong and the input list of
Theorem~\ref{thm:contradiction} is complete and final. All are
honest re-statements, not refutations; the Fork~A strategy and its
core constructions are unchanged and were independently confirmed
sound. An over-claim in an earlier Phase-1 draft --- a clean
``standard-isotypic'' landing theorem --- had already been caught and
withdrawn. Recommended next: a further independent re-validation of
this twice-revised document, then Phase~3 (Lean).

\end{document}
